













%%%%%%%%%%%%%%%%%%%%%%%%%%%%%%%%%%%%%%%%%%%%%%%%%%%%%%%%%%%%%%%%%%%%%%



\section{Appendix A: Proof of Theorem~\ref{thm:UsgnH-Ustar-MUstar-general}}
\label{sec:proof-thm:UsgnH-Ustar-MUstar-general}

%In the rest of this chapter, we establish the results Theorem~\ref{thm:UsgnH-Ustar-MUstar-general} and Corollary~\ref{cor:entrywise-error-general}.
To simplify notation, we assume throughout the proof that $\lambda_r^{\star} > 0$, namely,
%
\begin{align}
	|\lambda_1^{\star}| \geq  \cdots \geq | \lambda_{r-1}^{\star} | \geq \lambda_r^{\star} > 0.
	\label{eq:assumption-lambda-r-positive}
\end{align}
%Now we turn to the proof of Theorem~\ref{thm:UsgnH-Ustar-MUstar-general}.
The challenge of the proof arises due to the complicated statistical dependency between $\bm{M}$ and $\bm{U}$, and the leave-one-out analysis paves a plausible path to decouple the dependency.






\subsection{Construction of leave-one-out auxiliary estimates}
\label{sec:construction-LOO-general}

As elucidated in the rank-1 matrix denoising example in Section~\ref{sec:rank-1-denoising-LOO}, the key to enabling fine-grained analysis is to seek assistance from  a collection of leave-one-out estimates.  Akin to Section~\ref{sec:LOO-estimates-denoising}, for each $1\leq l\leq n$, we construct  two auxiliary matrices  $\bm{M}^{(l)}$ and $\bm{E}^{(l)}=\big[E^{(l)}_{i,j} \big]_{1\leq i,j\leq n}$ as follows:
%
\begin{align}
\bm{M}^{(l)}\coloneqq\bm{M}^{\star}+\bm{E}^{(l)},
\quad
E_{i,j}^{(l)}  \coloneqq
\begin{cases}
E_{i,j}, & \text{if }i\neq l\text{ and }j\neq l,\\
0, & \text{else},
\end{cases}
\label{eq:Ml-construction-general}
\end{align}
%
which are generated by simply discarding all random noise incurred in the $l$-th column/row of the data matrix.
In addition, let $\lambda_1^{(l)},\cdots,\lambda_n^{(l)}$ be the eigenvalues of $\bm{M}^{(l)}$ sorted by
%
\begin{align}
	\big|\lambda_1^{(l)}\big| \geq \big|\lambda_2^{(l)}\big| \geq \cdots \geq \big| \lambda_{n}^{(l)} \big| ,
\end{align}
%
and denote by $\bm{u}_{i}^{(l)}$ the eigenvector of $\bm{M}^{(l)}$  associated with $\lambda_i^{(l)}$. The leave-one-out spectral estimates $\bm{U}^{(l)}$ and $\bm{\Lambda}^{(l)}$ are, therefore, given by
%
\begin{align}
	\bm{U}^{(l)}\coloneqq\big[\bm{u}_{1}^{(l)},\cdots,\bm{u}_{r}^{(l)}\big]\in\mathbb{R}^{n\times r};
	\quad
	\bm{\Lambda}^{(l)}\coloneqq\mathsf{diag}\big(\big[\lambda_{1}^{(l)},\cdots,\lambda_{r}^{(l)}\big]\big).	
\end{align}
%

We emphasize again that the main advantage of introducing the leave-one-out estimate $\bm{U}^{(l)}$ stems from its statistical independence from the $l$-th row of $\bm{M}$, which substantially simplifies the analysis for the $l$-th row of the estimate.
In principle, our analysis employs the leave-one-out estimates
to help decouple delicate statistical dependency  in a row-by-row fashion. Another crucial aspect of the analysis lies in the exploitation of the proximity of all these auxiliary estimates, a feature that is enabled by the ``stability'' of the spectral method.





\subsection{Preliminary facts}
\label{sec:preliminary-facts-Linf-general}

%Before embarking on  the $\ell_{\infty}/\ell_{2,\infty}$ analysis, we gather a couple of useful facts immediately available from the $\ell_2$
%perturbation theory and the matrix tail bounds. In what follows, $\bm{\Theta}$ (resp.~$\bm{\Theta}^{(l)}$) denotes a diagonal matrix whose diagonal entries are the principal angles between $\bm{U}$ (resp.~$\bm{U}^{(l)}$) and $\bm{U}^{\star}$.
%%
%\begin{lemma}
%\label{lem:preliminary-result-general-iid}
%Consider the setting in Section~\ref{sec:setup-general-theory-independent}.
%There is some constant $c_2>0$ such that with probability at least $1-O(n^{-7})$,
%%
%\begin{subequations}
%\label{eq:L2-perturbation-summary-general}
%%
%\begin{align}
%	\|\bm{E} \|  &\leq c_2\sigma \sqrt{n}  \quad  & \big\|\bm{E}^{(l)} \big\| &\leq c_2 \sigma \sqrt{n}  &&  \label{eq:noise-size-general} \\
%  \mathsf{dist}( \bm{U} , \bm{U}^{\star} )  &\leq \frac{2c_2\sigma\sqrt{n}}{\lambda_r^{\star}}  ~~
% & \mathsf{dist}\big( \bm{U}^{(l)} , \bm{U}^{\star} \big)  &\leq \frac{2c_2\sigma\sqrt{n}}{\lambda_r^{\star}} && \label{eq:dist-U-Ustar-general} \\
%	  \big\|\sin \bm{\Theta} \big\|  &\leq \frac{c_2\sigma\sqrt{2n}}{\lambda_r^{\star}}  ~~
%	& \big\|\sin \bm{\Theta}^{(l)} \big\|  &\leq \frac{c_2\sigma\sqrt{2n}}{\lambda_r^{\star}} && \label{eq:distp-U-Ustar-general} \\
%	\max_{1\leq j\leq r} | \lambda_j  |  &\geq \lambda_r^{\star} - c_2 \sigma \sqrt{n} &  \max_{1\leq j\leq r} | \lambda^{(l)}_j  |  &\geq \lambda_r^{\star} - c_2\sigma \sqrt{n} &&  \label{eq:lambda-size-large-general}  \\
%	\max_{j: j> r} | \lambda_j   |  &\leq c_2 \sigma \sqrt{n} & \max_{j: j> r}| \lambda_j^{(l)} |  &\leq c_2\sigma \sqrt{n}  \label{eq:lambda-size-small-general}
%\end{align}
%%
%%
%hold simultaneously for all $1\leq l\leq n$, provided that $c_2\sigma \sqrt{n}\leq (1-1/\sqrt{2})\lambda_r^{\star}$.
%Moreover, for any fixed matrix $\bm{A}\in \mathbb{R}^{n\times d}$ with $d\leq n$, one has
%%
%\begin{align}
%	\|\bm{E} \bm{A} \|_{2,\infty} \leq
%	4\sigma\sqrt{\log n}\,\|\bm{A}\|_{\mathrm{F}}+(6B\log n)\|\bm{A}\|_{2,\infty}
%	\label{eq:E-concentration-inf2-general}
%\end{align}
%\end{subequations}
%%
%with probability exceeding $1-2n^{-5}$.
%\end{lemma}
%%
%\begin{remark}
%In view of \eqref{eq:E-concentration-inf2-general}, with probability at least $1-2n^{-5}$,
%%
%\begin{align}
%	& \|\bm{E}\bm{U}^{\star}\|_{2,\infty} \leq 4\sigma\sqrt{\log n}\,\|\bm{U}^{\star}\|_{\mathrm{F}}+(6B\log n)\|\bm{U}^{\star}\|_{2,\infty}
%		\nonumber \\
% 	&\qquad = 4\sigma\sqrt{r\log n}+6B\sqrt{\frac{\mu r\log^{2}n}{n}}
%	\leq (4+6c_{\mathsf{b}} ) \sigma \sqrt{r\log n},
%	\label{eq:EU-2-inf-bound-general}
%\end{align}
%%
%which relies on the definition \eqref{eq:defn-Ustar-incoherence} and \eqref{eq:assumption-B-sigma} in Assumption~\ref{assumption-noise-general}.  As a result,
%%
%\begin{align}
%\big\|\bm{M}\bm{U}^{\star}\big\|_{2,\infty} & \leq\big\|\bm{M}^{\star}\bm{U}^{\star}\big\|_{2,\infty}+\big\|\bm{E}\bm{U}^{\star}\big\|_{2,\infty} \notag\\
% % \leq\big\|\bm{U}^{\star}\big\|_{2,\infty}\|\bm{\Lambda}^{\star}\|+\big\|\bm{E}\bm{U}^{\star}\big\|_{2,\infty} \notag\\
% & \leq \sqrt{\frac{\mu r}{n}}\big|\lambda_{1}^{\star}\big|+(4+6c_{\mathsf{b}} )\sigma\sqrt{r\log n},
%	\label{eq:MU-2-inf-bound-general}
%\end{align}
%%
%which follows from the fact $\bm{M}^{\star}\bm{U}^{\star}=\bm{U}^{\star}\bm{\Lambda}^{\star}$ (so that $\|\bm{M}^{\star}\bm{U}^{\star}\|_{2,\infty}\leq \big\|\bm{U}^{\star}\big\|_{2,\infty}\|\bm{\Lambda}^{\star}\| = \sqrt{\mu r/n} \,|\lambda_1^{\star}|$).
%\end{remark}
%
%
%Furthermore, we find it helpful to introduce the following  matrices
%%
%\begin{equation}
%	\bm{H} \coloneqq \bm{U}^{\top}\bm{U}^{\star}
%	\qquad\text{and}\qquad
%	\bm{H}^{(l)} \coloneqq \bm{U}^{(l)\top}\bm{U}^{\star},
%	\label{eq:defn-H-Hl-general}
%\end{equation}
%%
%which turn out to be close to being orthonormal.
%%
%Our results concerning $\bm{H}$ and $\bm{H}^{(l)}$ are formalized as follows.
%%
%\begin{lemma}
%\label{lem:H-property-summary-general}
%Instate the assumptions of  Lemma~\ref{lem:preliminary-result-general-iid}. With probability at least $1-O(n^{-7})$, the following holds:
%%
%\begin{subequations}  \label{eq:H-property-summary-general}
%%
%\begin{align}
%	\|\bm{H}^{-1}\|  &\leq 2  \quad  & \big\|\big(\bm{H}^{(l)} \big) ^{-1} \big\| &\leq 2  &&  \label{eq:H-inv-norm-bound-general} \\
%  \big\|\bm{H}-\mathsf{sgn}(\bm{H})\big\|  &\leq  \frac{2c_{2}^{2}n \sigma^{2}}{(\lambda_{r}^{\star})^{2}}  ~~
%	&      \big\|\bm{H}^{(l)}-\mathsf{sgn}(\bm{H}^{(l)})\big\|  &  \leq  \frac{2c_{2}^{2}n\sigma^{2} }{(\lambda_{r}^{\star})^{2}}   &&
%	\label{eq:H-sgnH-dist-general}
%\end{align}
%%
%\end{subequations}
%%
%simultaneously for all $1\leq l\leq n$ and, in addition,
%%
%\begin{align}
%  \|\bm{U}\bm{H}-\bm{U}^{\star}\|_{\mathrm{F}}  \leq\frac{2c_{2}\sigma \sqrt{rn}}{\lambda_{r}^{\star}} .
%\end{align}
%%
%
%
%\end{lemma}
%
%\begin{remark}\label{remark:AH-norm-bound}
%As a consequence of \eqref{eq:H-property-summary-general} and Proposition~\ref{prop:unitary_norm_relation},  one has
%%
%\begin{subequations}  \label{eq:A-2inf-H-Hl-general}
%\begin{align}
%	\vertiii{\bm{A}} &= \vertiiibig{ \bm{A}\bm{H} \bm{H} ^{-1}}
%	\leq \vertiii{ \bm{A}\bm{H}} \, \big\| \bm{H} ^{-1}\big\|
%	\leq 2 \vertiii{ \bm{A}\bm{H} }
%	\label{eq:A-2inf-H-general} \\
%	\vertiii{\bm{A}}
%	& \leq \vertiiibig{ \bm{A}\bm{H}^{(l)}} \, \big\| \big( \bm{H}^{(l)}  \big) ^{-1}\big\|
%  	\leq 2 \vertiiibig{ \bm{A}\bm{H}^{(l)} }
%	\label{eq:A-2inf-Hl-general}
%\end{align}
%\end{subequations}
%%
%for any matrix $\bm{A}$ and any unitarily invariant norm $\vertiii{\cdot}$ (e.g., the Frobenius norm  and the $\ell_{2,\infty}$ norm $\|\cdot\|_{2,\infty}$).
%\end{remark}
%
%
%
%The proofs for the above results  are postponed to Section~\ref{sec:proof-auxiliary-lemmas-general}.
%
%
Before embarking on the $\ell_{\infty}$ and $\ell_{2,\infty}$ analyses,
we gather a couple of useful facts, whose proofs are postponed to
Section~\ref{sec:proof-auxiliary-lemmas-general}. In what follows,
$\bm{\Theta}$ (resp.~$\bm{\Theta}^{(l)}$) denotes a diagonal matrix
whose diagonal entries are the principal angles between $\bm{U}$
(resp.~$\bm{U}^{(l)}$) and $\bm{U}^{\star}$. In addition, we find
it helpful to introduce the following matrices
\begin{equation}
\bm{H}\coloneqq\bm{U}^{\top}\bm{U}^{\star}\qquad\text{and}\qquad\bm{H}^{(l)}\coloneqq\bm{U}^{(l)\top}\bm{U}^{\star},\label{eq:defn-H-Hl-general}
\end{equation}
which turn out to be close to being orthonormal.

The first set of results follows from the statistical nature of the
perturbation matrix $\bm{E}$ (cf.~Assumption~\ref{assumption-noise-general}),
which is immediately available from the matrix tail bounds. 
%
\begin{lemma}\label{lemma:general-noise-bound}
Consider the setting in Section~\ref{sec:setup-general-theory-independent}.
There is some constant $c_{2}>0$ such that with probability at least
$1-O(n^{-7})$,
\begin{equation}
\max_{l}\big\|\bm{E}^{(l)}\big\|\leq\|\bm{E}\|\leq c_{2}\sigma\sqrt{n}.\label{eq:noise-size-general}
\end{equation}
Moreover, for any fixed matrix $\bm{A}\in\mathbb{R}^{n\times d}$
with $d\leq n$, one has
\begin{align}
\|\bm{E}\bm{A}\|_{2,\infty}\leq4\sigma\sqrt{\log n}\,\|\bm{A}\|_{\mathrm{F}}+(6B\log n)\|\bm{A}\|_{2,\infty}.\label{eq:E-concentration-inf2-general}
\end{align}
 \end{lemma}
 %
%with probability at least $1-2n^{-5}$.

\begin{remark} In view of \eqref{eq:E-concentration-inf2-general},
with probability at least $1-2n^{-5}$,
\begin{align}
 & \|\bm{E}\bm{U}^{\star}\|_{2,\infty}\leq4\sigma\sqrt{\log n}\,\|\bm{U}^{\star}\|_{\mathrm{F}}+(6B\log n)\|\bm{U}^{\star}\|_{2,\infty}\nonumber \\
 & \qquad=4\sigma\sqrt{r\log n}+6B\sqrt{\frac{\mu r\log^{2}n}{n}}
 = (4+6c_{\mathsf{b}})\sigma\sqrt{r\log n},\label{eq:EU-2-inf-bound-general}
\end{align}
which relies on the definition \eqref{eq:defn-Ustar-incoherence}
and the definition of $c_{\mathsf{b}}$ in \eqref{eq:assumption-B-sigma}.
As a result,
\begin{align}
\big\|\bm{M}\bm{U}^{\star}\big\|_{2,\infty} & \leq\big\|\bm{M}^{\star}\bm{U}^{\star}\big\|_{2,\infty}+\big\|\bm{E}\bm{U}^{\star}\big\|_{2,\infty}\nonumber \\
&\leq \sqrt{\frac{\mu r}{n}}|\lambda_1^{\star}| + (4+6c_{\mathsf{b}})\sigma\sqrt{r\log n}, \label{eq:MU-2-inf-bound-general}
\end{align}
which follows from the fact $\bm{M}^{\star}\bm{U}^{\star}=\bm{U}^{\star}\bm{\Lambda}^{\star}$
(so that $\|\bm{M}^{\star}\bm{U}^{\star}\|_{2,\infty}\leq\big\|\bm{U}^{\star}\big\|_{2,\infty}\|\bm{\Lambda}^{\star}\|=\sqrt{\mu r/n}\,|\lambda_{1}^{\star}|$).
\end{remark}

With the size of the perturbations (i.e., $\|\bm{E}^{(l)}\|$ and
$\|\bm{E}\|$) under control, the $\ell_{2}$ perturbation theory
established in Chapter~\ref{cha:matrix-perturbation} leads to the following set of conclusions.
\begin{lemma}\label{lem:preliminary-result-general-iid} Suppose
that $c_{2}\sigma\sqrt{n}\leq(1-1/\sqrt{2})\lambda_{r}^{\star}$,
where $c_{2}$ is the same constant as in Lemma~\ref{lemma:general-noise-bound}.
Then with probability at least $1-O(n^{-7})$, one has \begin{subequations}
\label{eq:L2-perturbation-summary-general}
\begin{align}
\mathsf{dist}(\bm{U},\bm{U}^{\star}) & \leq\frac{2c_{2}\sigma\sqrt{n}}{\lambda_{r}^{\star}},~~ & \mathsf{dist}\big(\bm{U}^{(l)},\bm{U}^{\star}\big) & \leq\frac{2c_{2}\sigma\sqrt{n}}{\lambda_{r}^{\star}}, \label{eq:dist-U-Ustar-general}\\
\big\|\sin\bm{\Theta}\big\| & \leq\frac{c_{2}\sigma\sqrt{2n}}{\lambda_{r}^{\star}}, ~~ & \big\|\sin\bm{\Theta}^{(l)}\big\| & \leq\frac{c_{2}\sigma\sqrt{2n}}{\lambda_{r}^{\star}}, \label{eq:distp-U-Ustar-general}\\
\max_{1\leq j\leq r}|\lambda_{j}| & \geq\lambda_{r}^{\star}-c_{2}\sigma\sqrt{n}, & \max_{1\leq j\leq r}|\lambda_{j}^{(l)}| & \geq\lambda_{r}^{\star}-c_{2}\sigma\sqrt{n}, \label{eq:lambda-size-large-general}\\
\max_{j:j>r}|\lambda_{j}| & \leq c_{2}\sigma\sqrt{n}, & \max_{j:j>r}|\lambda_{j}^{(l)}| & \leq c_{2}\sigma\sqrt{n}\label{eq:lambda-size-small-general}
\end{align}
hold simultaneously for all $1\leq l\leq n$. In addition,
\begin{align}
\|\bm{U}\bm{H}-\bm{U}^{\star}\|_{\mathrm{F}}\leq\frac{2c_{2}\sigma\sqrt{rn}}{\lambda_{r}^{\star}}.\label{eq:dist-under-H-general}
\end{align}
 \end{subequations} %
with probability exceeding $1-2n^{-5}$. \end{lemma}


\begin{remark}\label{remark:perturbation-size-eigen-gap}
Lemmas~\ref{lemma:general-noise-bound} and \ref{lem:preliminary-result-general-iid} allow us to bound the eigengap and perturbation size as follows
%
\begin{subequations}\label{eq:lambda-rl-lambda-r1l-M-Ml-LB-step}
\begin{align}
	\big| \lambda_{r}^{(l)} \big| - \big|\lambda_{r+1}^{(l)}\big|
	& \geq \lambda_{r}^{\star}-  c_2 \sigma \sqrt{n} - c_2 \sigma \sqrt{n}
	\geq \lambda_{r}^{\star}/2,
	\label{eq:lambda-rl-lambda-r1l-LB-step}  \\
	\|\bm{M}-\bm{M}^{(l)}\| & \leq\|\bm{M}-\bm{M}^{\star}\|+\|\bm{M}^{\star}-\bm{M}^{(l)}\|=\|\bm{E}\|+\|\bm{E}^{(l)}\|  \notag\\
	& \leq2c_{2}\sigma\sqrt{n}\leq(1-1/\sqrt{2}) \big( \big| \lambda_{r}^{(l)} \big| - \big|\lambda_{r+1}^{(l)}\big| \big),
	\label{eq:M-Ml-size-step}
\end{align}
%
\end{subequations}
%
which are valid as long as $20 c_2 \sigma \sqrt{n} \leq \lambda_r^{\star}$. These will prove useful when bounding the approximation error of $\bm{U}$ using $\bm{U}^{(l)}$.
\end{remark}

Another collection of results is concerned with $\bm{H}$ and $\bm{H}^{(l)}$.
%
\begin{lemma} \label{lem:H-property-summary-general} Suppose that the
assumptions of Lemma~\ref{lem:preliminary-result-general-iid} hold. With
probability at least $1-O(n^{-7})$, \begin{subequations}
\label{eq:H-property-summary-general}
\begin{align}
\|\bm{H}^{-1}\| & \leq2,\quad & \big\|\big(\bm{H}^{(l)}\big)^{-1}\big\| & \leq2,\label{eq:H-inv-norm-bound-general}\\
\big\|\bm{H}-\mathsf{sgn}(\bm{H})\big\| & \leq\frac{2c_{2}^{2}n\sigma^{2}}{(\lambda_{r}^{\star})^{2}},~~ & \big\|\bm{H}^{(l)}-\mathsf{sgn}(\bm{H}^{(l)})\big\| & \leq\frac{2c_{2}^{2}n\sigma^{2}}{(\lambda_{r}^{\star})^{2}}\label{eq:H-sgnH-dist-general}
\end{align}
\end{subequations} hold simultaneously for all $1\leq l\leq n$.

\end{lemma}

\begin{remark}\label{remark:AH-norm-bound} As a consequence of \eqref{eq:H-property-summary-general}
and Proposition~\ref{prop:unitary_norm_relation}, one has \begin{subequations}
\label{eq:A-2inf-H-Hl-general}
\begin{align}
\vertiii{\bm{A}} & =\vertiiibig{\bm{A}\bm{H}\bm{H}^{-1}}\leq\vertiii{\bm{A}\bm{H}}\,\big\|\bm{H}^{-1}\big\|\leq2\vertiii{\bm{A}\bm{H}}, \label{eq:A-2inf-H-general}\\
\vertiii{\bm{A}} & \leq\vertiiibig{\bm{A}\bm{H}^{(l)}}\,\big\|\big(\bm{H}^{(l)}\big)^{-1}\big\|\leq2\vertiiibig{\bm{A}\bm{H}^{(l)}}\label{eq:A-2inf-Hl-general}
\end{align}
\end{subequations} for any matrix $\bm{A}$. Here, $\vertiii{\bm{A}}$ could either be the Frobenius norm or the $\ell_{2,\infty}$
norm $\|\cdot\|_{2,\infty}$. \end{remark}




\subsection{Leave-one-out analysis}

Now we move on to the main part of the analysis, which is further decomposed into four steps. 

\subsubsection{Step 1: decomposing the $\ell_{2,\infty}$ estimation error of $\bm{U}$}




By virtue of the proximity of  $\bm{H}$ and $\mathsf{sgn}(\bm{H})$ unveiled in Lemma~\ref{lem:H-property-summary-general}, we are allowed to employ $\bm{U}\bm{H}$ as a surrogate for $\bm{U}\mathsf{sgn}(\bm{H})$, which is more convenient to work with. As it turns out, the discrepancy between $\bm{U}\bm{H}$ and the first-order approximation $\bm{M}\bm{U}^{\star}(\bm{\Lambda}^{\star})^{-1}$, and the discrepancy between $\bm{U}\bm{H}$ and the truth,  can be bounded by three important terms, as asserted below. The proof is built upon elementary  algebra and basic $\ell_2$ perturbation bounds in Section~\ref{sec:preliminary-facts-Linf-general}, and is deferred to Section~\ref{sec:proof-auxiliary-lemmas-general}.




\begin{lemma}
\label{lem:UH-MUstar-diff-decompose}
%
Suppose that $2c_2\sigma\sqrt{n}\leq \lambda_r^{\star}$ for some sufficiently large constant $c_2>0$. Then with probability at least $1-O(n^{-7})$, one has
%
\begin{subequations}
	\label{eq:UH-MUstar-Ustar-diff-2inf-general}
	\begin{align}
		\big\|\bm{U}\bm{H}-\bm{M}\bm{U}^{\star}\big(\bm{\Lambda}^{\star}\big)^{-1}\big\|_{2,\infty} & \leq\mathcal{E}_{1}+\mathcal{E}_{2},
		\label{eq:UH-MUstarLambdastar-diff-2inf-general}\\
		\big\|\bm{U}\bm{H}-\bm{U}^{\star}\big\|_{2,\infty} & \leq \mathcal{E}_{1}+\mathcal{E}_{2}+ \mathcal{E}_3,
		\label{eq:UH-Ustar-diff-2inf-general}
	\end{align}
\end{subequations}
%
where
%
%\begin{subequations}
%\begin{align*}
%	\label{eq:UB-UH-MULambda-UB-lemma}
%
%\begin{array}{cc}
$\mathcal{E}_{1}\coloneqq\frac{2\|\bm{M}(\bm{U}\bm{H}-\bm{U}^{\star})\|_{2,\infty}}{\lambda_{r}^{\star}}$, $\mathcal{E}_{2}\coloneqq\frac{4\|\bm{M}\bm{U}^{\star}\|_{2,\infty}\|\bm{E}\|}{(\lambda_{r}^{\star})^{2}}$, and $\mathcal{E}_{3}\coloneqq\frac{\|\bm{E}\bm{U}^{\star}\|_{2,\infty}}{\lambda_{r}^{\star}}$.
%\end{array}
%
%\end{align*}
%\end{subequations}
%
\end{lemma}
%




Lemma~\ref{lem:UH-MUstar-diff-decompose} leaves us with three important terms to deal with. The term $\mathcal{E}_1$ is most complicated as it involves the product of two random matrices, whereas $\mathcal{E}_2$ and $\mathcal{E}_3$ can be controlled straightforwardly through our preliminary facts in Section~\ref{sec:preliminary-facts-Linf-general}.   Specifically,   the term $\mathcal{E}_2$ can be bounded by combining \eqref{eq:MU-2-inf-bound-general} with the bound \eqref{eq:noise-size-general} on $\bm{E}$ to obtain
%
% \begin{align*}
% \big\|\bm{M}\bm{U}^{\star}\big\|_{2,\infty} & \leq\big\|\bm{M}^{\star}\bm{U}^{\star}\big\|_{2,\infty}+\big\|\bm{E}\bm{U}^{\star}\big\|_{2,\infty}
%   \leq\big\|\bm{U}^{\star}\big\|_{2,\infty}\|\bm{\Lambda}^{\star}\|+\big\|\bm{E}\bm{U}^{\star}\big\|_{2,\infty}\\
%  & \leq \sqrt{\frac{\mu r}{n}}\big|\lambda_{1}^{\star}\big|+(4+6c_{1})\sigma\sqrt{r\log n},
% \end{align*}
%
% where the first identity makes use of the fact $\bm{M}^{\star}\bm{U}^{\star}=\bm{U}^{\star}\bm{\Lambda}^{\star}$ (so that $\|\bm{M}^{\star}\bm{U}^{\star}\|_{2,\infty}\leq \big\|\bm{U}^{\star}\big\|_{2,\infty}\|\bm{\Lambda}^{\star}\|$), and the last line relies on \eqref{eq:EU-2-inf-bound-general}. This combined with
% the bound \eqref{eq:L2-perturbation-summary-general} on $\|\bm{E}\|$ yields
%
\begin{align}
	% \frac{\big\|\bm{M}\bm{U}^{\star}\big\|_{2,\infty}\|\bm{E}\|}{(\lambda_{r}^{\star})^{2}}
	\mathcal{E}_2 & \leq \frac{4c_{2}\kappa\sigma\sqrt{\mu r}}{\lambda_{r}^{\star}}+\frac{4c_{2}(4+6c_{\mathsf{b}})\sigma^{2}\sqrt{rn\log n}}{(\lambda_{r}^{\star})^{2}}.
	\label{eq:MUstar-E-norm-lambda2-bound}
\end{align}
%
Regarding the term $\mathcal{E}_3$, the inequality \eqref{eq:EU-2-inf-bound-general} readily gives
%
\begin{align}
	%\big\|\bm{E}\bm{U}^{\star}\big(\bm{\Lambda}^{\star}\big)^{-1}\big\|_{2,\infty}\leq\frac{\big\|\bm{E}\bm{U}^{\star}\big\|_{2,\infty}}{\lambda_{r}^{\star}}
	\mathcal{E}_3 \leq\frac{(4+6c_{\mathsf{b}} )\sigma\sqrt{r\log n}}{\lambda_{r}^{\star}}.
	\label{eq:E3-bound-general}
\end{align}
%




Turning to controlling the remaining term $\mathcal{E}_{1}$, a closer inspection, however, reveals substantial challenges, due to the complicated statistical dependency between $\bm{M}$ and $\bm{U}$. To further complicate matters, the term $\mathcal{E}_1$---as we shall demonstrate momentarily---depend on some intrinsic properties of interest about $\bm{U}$ (e.g., $\|\bm{U}\bm{H}-\bm{U}^{\star}\|_{2,\infty}$),  which might lead to circular reasoning if not handled properly. In order to circumvent this issue, we intend to establish the following relation
%
\begin{equation}
	\mathcal{E}_{1}\leq\mathcal{E}_{1,1}+\rho_{1}\big\|\bm{U}\bm{H}-\bm{U}^{\star}\big\|_{2,\infty}
	\label{eq:E1-UB-E11-rho1}
\end{equation}
%
for some quantity $\mathcal{E}_{1,1}>0$ that does not involve $\|\bm{U}\bm{H}-\bm{U}^{\star} \|_{2,\infty}$  as well as some contraction factor $0<\rho_{1}\leq1/2$. Assuming the relation~\eqref{eq:E1-UB-E11-rho1} holds for the moment, we have the following useful claim (the proof is straightforward and again postponed to Section~\ref{sec:proof-auxiliary-lemmas-general}).
%
\begin{lemma}
\label{lem:UH-sequence-bounds-E123}
If Conditions \eqref{eq:UH-MUstar-Ustar-diff-2inf-general} and \eqref{eq:E1-UB-E11-rho1} hold with $0<\rho_{1}\leq1/2$, then we have
%
\begin{subequations}\label{eq:UH-sequence-bounds-E123}
\begin{align}
\big\|\bm{U}\bm{H}-\bm{U}^{\star}\big\|_{2,\infty} & \leq2\big(\mathcal{E}_{1,1}+\mathcal{E}_{2}+\mathcal{E}_{3}\big), \label{eq:UH-Ustar-E11-E2-E3}\\
\big\|\bm{U}\bm{H}-\bm{M}\bm{U}^{\star}\big(\bm{\Lambda}^{\star}\big)^{-1}\big\|_{2,\infty} & \leq2\big(\mathcal{E}_{1,1}+\mathcal{E}_{2}+\rho_{1}\mathcal{E}_{3}\big),
\label{eq:UH-MUstar-Lambdastar-E11-E2-E3}
\end{align}
\begin{align}
	& \big\|\bm{U}\mathsf{sgn}(\bm{H})-\bm{U}^{\star}\big\|_{2,\infty}  \leq4\big(\mathcal{E}_{1,1}+\mathcal{E}_{2}+\mathcal{E}_{3}\big)+\frac{4c_{2}^{2}\sigma^{2}\sqrt{\mu rn}}{(\lambda_{r}^{\star})^{2}}, \label{eq:UsgnH-Ustar-Lambdastar-E11-E2-E3} \\
	& \big\|\bm{U}\mathsf{sgn}(\bm{H})-\bm{M}\bm{U}^{\star}\big(\bm{\Lambda}^{\star}\big)^{-1}\big\|_{2,\infty} \nonumber \\
& \qquad \leq3\mathcal{E}_{1,1}+3\mathcal{E}_{2}+\Big(2\rho_{1}+\frac{8c_{2}^{2}\sigma^{2}n}{(\lambda_{r}^{\star})^{2}}\Big)\mathcal{E}_{3}
 +\frac{4c_{2}^{2}\sigma^{2}\sqrt{\mu rn}}{(\lambda_{r}^{\star})^{2}}.\label{eq:UsgnH-MUstar-Lambdastar-E11-E2-E3}
\end{align}
\end{subequations}
%
\end{lemma}
%
\begin{remark}
When $\mathcal{E}_3$ is the dominant term, the  bound   \eqref{eq:UH-MUstar-Lambdastar-E11-E2-E3}  might be stronger than \eqref{eq:UH-Ustar-E11-E2-E3} if $\rho_1$ is  small. % This subtle improvement proves beneficial for some applications like community detection.
\end{remark}


% Let us take a closer inspection of the three terms $\mathcal{E}_{1}$, $\mathcal{E}_{2}$ and $\mathcal{E}_3$.




%Suppose for the moment that there exist some quantities $\mathcal{E}_{1,1}>0$ and $0<\rho_{1}\leq1/2$ obeying








With this lemma in mind, everything boils down to (i) establishing the relation \eqref{eq:E1-UB-E11-rho1} and (ii) deriving a tight bound on $\mathcal{E}_{1,1}$, which form the main content of the rest of the proof. In light of the triangle inequality
%
\begin{align*}
	\big\|\bm{M}\big(\bm{U}\bm{H}-\bm{U}^{\star}\big)\big\|_{2,\infty}
	\leq
	\big\|\bm{E} \big(\bm{U}\bm{H}-\bm{U}^{\star}\big)\big\|_{2,\infty}
	+ \big\|\bm{M}^{\star}\big(\bm{U}\bm{H}-\bm{U}^{\star}\big)\big\|_{2,\infty},
\end{align*}
%
we dedicate the next two steps to bounding $\|\bm{E} (\bm{U}\bm{H}-\bm{U}^{\star} ) \|_{2,\infty}$ and $\|\bm{M}^{\star} (\bm{U}\bm{H}-\bm{U}^{\star}) \|_{2,\infty}$ respectively.





\subsubsection{Step 2: bounding $\|\bm{E} (\bm{U}\bm{H}-\bm{U}^{\star}) \|_{2,\infty}$ via leave-one-out analysis}

To obtain tight row-wise control of $\bm{E}\big(\bm{U}\bm{H}-\bm{U}^{\star}\big)$, one needs to carefully decouple the statistical dependency between $\bm{E}$ and $\bm{U}$,  which is where the leave-one-out idea comes into play.



\paragraph{Step 2.1: a convenient decomposition.}

We start by invoking the triangle inequality to decompose the target quantity as follows
%
\begin{align}
 & \big\|\bm{E}\big(\bm{U}\bm{H}-\bm{U}^{\star}\big)\big\|_{2,\infty}=\max_{l}\big\|\bm{E}_{l,\cdot}\big(\bm{U}\bm{H}-\bm{U}^{\star}\big)\big\|_{2}\nonumber \\
 & \leq\max_{l}\Big\{\big\|\bm{E}_{l,\cdot}\big(\bm{U}^{(l)}\bm{H}^{(l)}-\bm{U}^{\star}\big)\big\|_{2}+\big\|\bm{E}_{l,\cdot}\big(\bm{U}\bm{H}-\bm{U}^{(l)}\bm{H}^{(l)}\big)\big\|_{2}\Big\}\nonumber \\
 & \leq\max_{l}\Big\{\big\|\bm{E}_{l,\cdot}\big(\bm{U}^{(l)}\bm{H}^{(l)}-\bm{U}^{\star}\big)\big\|_{2}+\|\bm{E}\| \big\|\bm{U}\bm{H}-\bm{U}^{(l)}\bm{H}^{(l)}\big\|_{\mathrm{F}}\Big\}.
	\label{eq:decomposition-E-UH-Ustar}
\end{align}
%
In words, when controlling the $l$-th row of $\bm{E}\big(\bm{U}\bm{H}-\bm{U}^{\star}\big)$, we attempt to employ $\bm{U}^{(l)}\bm{H}^{(l)}$ as a surrogate of $\bm{U}\bm{H}$.  The benefits to be harvested from this decomposition are:
%
\begin{itemize}
	\item The statistical independence between $\bm{E}_{l,\cdot}$ and $\bm{U}^{(l)}\bm{H}^{(l)}$ allows for convenient upper bounds on $\big\|\bm{E}_{l,\cdot}\big(\bm{U}^{(l)}\bm{H}^{(l)}-\bm{U}^{\star}\big)\big\|_{2}$;
	
	\item  $\bm{U}\bm{H}$ and  $\bm{U}^{(l)}\bm{H}^{(l)}$ are expected to be exceedingly close, so that the discrepancy incurred by replacing $\bm{U}\bm{H}$ with $\bm{U}^{(l)}\bm{H}^{(l)}$ is negligible.
\end{itemize}
%
In what follows, we flesh out the proof details.

%With regards to the second term in \eqref{eq:MUh-Ustar-UB12},





%Our starting point is the following decomposition
%%
%\begin{align}
%	& \big\|\bm{M}\big(\bm{U}\bm{H}-\bm{U}^{\star}\big) \big\|_{2,\infty}  \notag\\
%  %\leq\big\|\bm{M}\big(\bm{U}\bm{H}-\bm{U}^{\star}\big)\big\|_{2,\infty}\big\|\big(\bm{\Lambda}^{\star}\big)^{-1}\big\|\\
% %& \qquad
%	%=\frac{\big\|\bm{M}\big(\bm{U}\bm{H}-\bm{U}^{\star}\big)\big\|_{2,\infty}}{\lambda_{r}^{\star}}
%	& \quad \leq   \big\|\bm{M}^{\star}\big(\bm{U}\bm{H}-\bm{U}^{\star}\big)\big\|_{2,\infty} + \big\|\bm{E}\big(\bm{U}\bm{H}-\bm{U}^{\star}\big)\big\|_{2,\infty}.
%	\label{eq:MUh-Ustar-UB12}
%\end{align}
%%


%In what follows, we shall take care of each of these two terms separately.
%




\paragraph{Step 2.2: the proximity of $\bm{U}\bm{H}$ and $\bm{U}^{(l)}\bm{H}^{(l)}$.}
Given that
%
\begin{align}
	\big\|\bm{U}\bm{H}-\bm{U}^{(l)}\bm{H}^{(l)}\big\|_{\mathrm{F}}
	& =\big\|\bm{U}\bm{U}^{\top}\bm{U}^{\star}-\bm{U}^{(l)}\bm{U}^{(l)\top}\bm{U}^{\star}\big\|_{\mathrm{F}} \notag\\
 	& \leq\big\|\bm{U}\bm{U}^{\top}-\bm{U}^{(l)}\bm{U}^{(l)\top}\big\|_{\mathrm{F}} \big\|\bm{U}^{\star}\big\| \notag\\
 	& =\big\|\bm{U}\bm{U}^{\top}-\bm{U}^{(l)}\bm{U}^{(l)\top}\big\|_{\mathrm{F}},
	\label{eq:UU-UlUl-UB01}
\end{align}
%
it boils down to bounding $\|\bm{U}\bm{U}^{\top}-\bm{U}^{(l)}\bm{U}^{(l)\top}\|_{\mathrm{F}}$.
Under simple conditions on the eigengap and the perturbation size (see \eqref{eq:lambda-rl-lambda-r1l-M-Ml-LB-step}  in Remark~\ref{remark:perturbation-size-eigen-gap}),  the Davis-Kahan theorem (cf.~Corollary~\ref{cor:davis-kahan-conclusion-corollary}) yields
%
\begin{align}
	\big\|\bm{U}\bm{U}^{\top}-\bm{U}^{(l)}\bm{U}^{(l)\top}\big\|_{\mathrm{F}}
	& \leq\frac{2\big\|\big(\bm{M}-\bm{M}^{(l)}\big)\bm{U}^{(l)}\big\|_{\mathrm{F}} }{ \big| \lambda_{r}^{(l)} \big| - \big|\lambda_{r+1}^{(l)} \big|} \nonumber\\
	& \leq\frac{4\big\|\big(\bm{M}-\bm{M}^{(l)}\big)\bm{U}^{(l)}\big\|_{\mathrm{F}} }{\lambda_{r}^{\star}}.
	\label{eq:UU-UlUl-UB1}
\end{align}
%



It remains to develop an upper bound on $\| (\bm{M}-\bm{M}^{(l)} )\bm{U}^{(l)} \|$.
The way we construct $\bm{M}^{(l)}$ (see Section~\ref{sec:construction-LOO-general}) allows us to express
%
\[
	\big(\bm{M}-\bm{M}^{(l)}\big)\bm{U}^{(l)}=\bm{e}_{l}\bm{E}_{l,\cdot}\bm{U}^{(l)}+  \big(\bm{E}_{\cdot,l} - E_{l,l} \bm{e}_l \big)  \bm{e}_{l}^{\top}\bm{U}^{(l)}.
\]
%
This together with the triangle inequality and the fact \eqref{eq:A-2inf-H-Hl-general} gives
%
\begin{align*}
 & \big\|\big(\bm{M}-\bm{M}^{(l)}\big)\bm{U}^{(l)}\big\|_{\mathrm{F}}\leq\big\|\bm{E}_{l,\cdot}\bm{U}^{(l)}\big\|_{2}+\big\|\bm{E}_{\cdot,l} - E_{l,l} \bm{e}_l \big\|_{2}\,\big\|\bm{U}^{(l)}\big\|_{2,\infty}\\
%	& \quad\leq\big\|\bm{E}_{l,\cdot}\bm{U}^{(l)}\big\|_{2}+\big\|\bm{E}\big\|\cdot\big\|\bm{U}^{(l)} \bm{H}^{(l)}  \big\|_{2,\infty} \big\| \big(\bm{H}^{(l)}\big)^{-1} \big\| \\
 & \quad\leq\big\|\bm{E}_{l,\cdot}\bm{U}^{(l)}\big\|_{2}+2\big\|\bm{E}\big\|\, \big\|\bm{U}^{(l)}\bm{H}^{(l)}\big\|_{2,\infty}\\
 & \quad\leq\big\|\bm{E}_{l,\cdot}\bm{U}^{(l)}\big\|_{2}+2\big\|\bm{E}\big\|\, \big\|\bm{U}\bm{H}\big\|_{2,\infty}+2\big\|\bm{E}\big\|\,\big\|\bm{U}\bm{H}-\bm{U}^{(l)}\bm{H}^{(l)}\big\|_{\mathrm{F}}.
\end{align*}
%
%where the last line arises from the triangle inequality.
%
Substitution into \eqref{eq:UU-UlUl-UB01} and \eqref{eq:UU-UlUl-UB1} gives
%
\begin{align*}
 & \big\|\bm{U}\bm{H}-\bm{U}^{(l)}\bm{H}^{(l)}\big\|_{\mathrm{F}} \leq \big\|\bm{U}\bm{U}^{\top}-\bm{U}^{(l)}\bm{U}^{(l)\top}\big\|_{\mathrm{F}} \\
 & \quad\leq\frac{4\big\|\bm{E}_{l,\cdot}\bm{U}^{(l)}\big\|_{2}+8\big\|\bm{E}\big\|\,\big\|\bm{U}\bm{H}\big\|_{2,\infty}+8\big\|\bm{E}\big\|\, \big\|\bm{U}\bm{H}-\bm{U}^{(l)}\bm{H}^{(l)}\big\|_{\mathrm{F}}}{\lambda_{r}^{\star}}.
\end{align*}
%
As long as $\|\bm{E}\|/\lambda_r^{\star}\leq 1/16$,
%(which is guaranteed if $16c_2\sigma\sqrt{n}\leq \lambda_r^{\star}$),
 one can further rearrange terms to obtain
%
\begin{align}
\big\|\bm{U}\bm{H}-\bm{U}^{(l)}\bm{H}^{(l)}\big\|_{\mathrm{F}} & \leq\frac{8\big\|\bm{E}_{l,\cdot}\bm{U}^{(l)}\big\|_{2}+16\big\|\bm{E}\big\|\, \big\|\bm{U}\bm{H}\big\|_{2,\infty}}{\lambda_{r}^{\star}}.
	\label{eq:UH-UlHl-diff-UB10}
\end{align}
%
In addition, the fact  \eqref{eq:A-2inf-H-Hl-general} combined with the triangle inequality yields
%
\[
	\tfrac{1}{2}\big\|\bm{E}_{l,\cdot}\bm{U}^{(l)}\big\|_{2}\leq \big\|\bm{E}_{l,\cdot}\bm{U}^{(l)}\bm{H}^{(l)}\big\|_{2}
	\leq \big\|\bm{E}_{l,\cdot}\big(\bm{U}^{(l)}\bm{H}^{(l)}-\bm{U}^{\star}\big)\big\|_{2}+ \big\|\bm{E}_{l,\cdot}\bm{U}^{\star}\big\|_{2},
\]
%
which taken collectively with \eqref{eq:UH-UlHl-diff-UB10} reveals that
%\yxc{check whether to use $UH$ or $U$}
%
\begin{align}
 & \big\|\bm{U}\bm{H}-\bm{U}^{(l)}\bm{H}^{(l)}\big\|_{\mathrm{F}}
	\leq\frac{16\big\|\bm{E}_{l,\cdot}\big(\bm{U}^{(l)}\bm{H}^{(l)}-\bm{U}^{\star}\big)\big\|_{2}}{\lambda_{r}^{\star}} \nonumber \\
	&  \quad + \frac{ 16 \Big\{ \big\|\bm{E}_{l,\cdot}\bm{U}^{\star}\big\|_{2}  + \big\|\bm{E}\big\| \, \big\|\bm{U}\bm{H}-\bm{U}^{\star}\big\|_{2,\infty} +  \big\|\bm{E}\big\|\, \big\|\bm{U}^{\star}\big\|_{2,\infty} \Big\} }{\lambda_{r}^{\star}}.
	\label{eq:UH-UlHl-diff-UB123}
\end{align}
%










\paragraph{Step 2.3: bounding $\|\bm{E}_{l,\cdot} (\bm{U}^{(l)}\bm{H}^{(l)}-\bm{U}^{\star} ) \|_{2}$.}

%Another term that appears in  \eqref{eq:decomposition-E-UH-Ustar}
%
%is $\|\bm{E}_{l,\cdot} (\bm{U}^{(l)}\bm{H}^{(l)}-\bm{U}^{\star} ) \|_{2}$, which we look into in this step.
Recognizing that $\bm{E}_{l,\cdot}$ is statistically independent of $\bm{U}^{(l)}$ (since $\bm{U}^{(l)}$ is computed without using $\bm{E}_{l,\cdot}$),
we invoke Lemma~\ref{lemma:general-noise-bound} (more precisely, we use the proof of this lemma) to demonstrate that
with probability exceeding $1-2n^{-6}$,
%
\begin{align}
 & \big\|\bm{E}_{l,\cdot}\big(\bm{U}^{(l)}\bm{H}^{(l)}-\bm{U}^{\star}\big)\big\|_{2} \notag\\
	%\leq4\sqrt{v\log n}+6L\log n\nonumber \\
 & \leq4\sigma\sqrt{\log n} \, \|\bm{U}^{(l)}\bm{H}^{(l)}-\bm{U}^{\star}\|_{\mathrm{F}}
	+ (6B \log n) \|\bm{U}^{(l)}\bm{H}^{(l)}-\bm{U}^{\star}\|_{2,\infty} \nonumber \\
	& \leq4\sigma \sqrt{\log n} \, \|\bm{U}\bm{H}-\bm{U}^{\star}\|_{\mathrm{F}} + (6B \log n) \|\bm{U}\bm{H}-\bm{U}^{\star}\|_{2,\infty} \nonumber \\
	& \qquad+ ( 10B\log n) \|\bm{U}\bm{H}-\bm{U}^{(l)}\bm{H}^{(l)}\|_{\mathrm{F}}
 \label{eq:El-UH-Ustar-UB-Bern1}
\end{align}
%
holds simultaneously for all $1\leq l\leq n$, where the last line results from the triangle inequality and the fact $4\sigma\sqrt{\log n}+6B\log n\leq 10B\log n$.


% using Assumption~\ref{eq:assumption-B-sigma}


\paragraph{Step 2.4: combining the above bounds.}

The careful reader would immediately remark that the inequalities \eqref{eq:UH-UlHl-diff-UB123} and \eqref{eq:El-UH-Ustar-UB-Bern1} are convoluted, both of which  involve the terms  $\|\bm{E}_{l,\cdot} (\bm{U}^{(l)}\bm{H}^{(l)}-\bm{U}^{\star} ) \|_{2}$ and $\|\bm{U}\bm{H}-\bm{U}^{(l)}\bm{H}^{(l)}\|_{\mathrm{F}}$. Fortunately, one can substitute \eqref{eq:El-UH-Ustar-UB-Bern1} into \eqref{eq:UH-UlHl-diff-UB123} to produce a cleaner bound. By doing so and exploiting the condition $320B\log n \leq \lambda_r^{\star}$, we  rearrange terms to reach
%
\begin{align*}
	& \big\|\bm{U}\bm{H}-\bm{U}^{(l)}\bm{H}^{(l)}\big\|_{\mathrm{F}}
	\leq \frac{32\big\|\bm{E}_{l,\cdot}\bm{U}^{\star}\big\|_{2}+32\big\|\bm{E}\big\|\, \big\|\bm{U}^{\star}\big\|_{2,\infty}}{\lambda_{r}^{\star}}\nonumber \\
	& +\frac{128\sigma\sqrt{\log n}\big\|\bm{U}\bm{H}-\bm{U}^{\star}\big\|_{\mathrm{F}}+(32c_2\sigma \sqrt{n} +192B\log n)\big\|\bm{U}\bm{H}-\bm{U}^{\star}\big\|_{2,\infty}}{\lambda_{r}^{\star}},
	%\label{eq:UH-UlHl-diff-UB456}
\end{align*}
%
where we have also used the upper bound on $\|\bm{E}\|$ derived in \eqref{eq:L2-perturbation-summary-general}.
Meanwhile, plugging the above inequality into \eqref{eq:El-UH-Ustar-UB-Bern1} yields
%
\begin{align}
	& \big\|\bm{E}_{l,\cdot}\big(\bm{U}^{(l)}\bm{H}^{(l)}-\bm{U}^{\star}\big)\big\|_{2}
	\leq \frac{320B\log n}{\lambda_{r}^{\star}}\big(\big\|\bm{E}_{l,\cdot}\bm{U}^{\star}\big\|_{2}+\big\|\bm{E}\big\|\,\big\|\bm{U}^{\star}\big\|_{2,\infty}\big) \nonumber\\
	&\quad + 5\sigma\sqrt{\log n}\,\|\bm{U}\bm{H}-\bm{U}^{\star}\|_{\mathrm{F}}+(7B\log n)\|\bm{U}\bm{H}-\bm{U}^{\star}\|_{2,\infty}  ,
	\label{eq:El-UH-Ustar-UB-Bern123}
\end{align}
%
provided that $\max \{ \sigma \sqrt{n}, B\log n\} \leq c_3\lambda_r^{\star}$ for some  small constant $c_3$.



Substituting the preceding two bounds into \eqref{eq:decomposition-E-UH-Ustar} and combining terms reveal the existence of some constant $c_4>0$ such that
%
\begin{align}
	&  \big\|\bm{E}\big(\bm{U}\bm{H}-\bm{U}^{\star}\big)\big\|_{2,\infty} \notag\\
	& \qquad \leq \alpha_0 + \alpha_1 + c_4\big( \sigma  \sqrt{n} + B\log n \big)\|\bm{U}\bm{H}-\bm{U}^{\star}\|_{2,\infty}
	\label{eq:El-UH-Ustar-UB-Bern123}
\end{align}
%
provided that $\max \{ \sigma \sqrt{n}, B\log n\} \leq c_3\lambda_r^{\star}$ for some constant $c_3>0$ small enough, where
\begin{align}
\begin{array}{ll}
 & \alpha_{0}\coloneqq\frac{32c_{2}\sigma\sqrt{n}+320B\log n}{\lambda_{r}^{\star}}\big(\big\|\bm{E}\bm{U}^{\star}\big\|_{2,\infty}+\big\|\bm{E}\big\|\,\big\|\bm{U}^{\star}\big\|_{2,\infty}\big) , \\
 & \alpha_{1}\coloneqq6\sigma\sqrt{\log n}\,\|\bm{U}\bm{H}-\bm{U}^{\star}\|_{\mathrm{F}} .
\end{array}	
\label{eq:defn-alpha-0-1-general}
\end{align}
%


Before continuing,  note that we are already well-equipped to bound the above two quantities. First, $\alpha_0$ can be bounded by
%
\begin{align}
\alpha_{0} & \leq \frac{32c_{2}\sigma\sqrt{n}+320B\log n}{\lambda_{r}^{\star}}\nonumber \\
 & \qquad\cdot\Big(4\sigma\sqrt{r\log n}+6B\sqrt{\frac{\mu r}{n}}\log n+c_{2}\sigma\sqrt{\mu r}\Big),
\label{eq:defn-alpha-0-general}
\end{align}
%
where we have used the bounds concerning $\bm{E}$ from Lemma~\ref{lemma:general-noise-bound} as well as the definition \eqref{eq:defn-Ustar-incoherence}.
Regarding $\alpha_1$, it is seen from (\ref{eq:dist-under-H-general}) that
%
\begin{align}
\alpha_{1}
  \leq\frac{12c_{2}\sigma^{2}\sqrt{rn\log n}}{\lambda_{r}^{\star}}.
	\label{eq:alpha1-bound-1234}
\end{align}


%Substituting the inequalities \eqref{eq:UH-UlHl-diff-UB123} and \eqref{eq:El-UH-Ustar-UB-Bern1} into \eqref{eq:decomposition-E-UH-Ustar} with a little algebra yields
%%
%\begin{align}
% & \big\|\bm{E}(\bm{U}\bm{H}-\bm{U}^{\star})\big\|_{2,\infty}\nonumber \\
% & \quad\leq4\sigma\sqrt{\log n}\|\bm{U}\bm{H}-\bm{U}^{\star}\|_{\mathrm{F}}+(6B\log n)\|\bm{U}\bm{H}-\bm{U}^{\star}\|_{2,\infty}\nonumber \\
% & \qquad+\frac{4\big(\|\bm{E}\|+10B\log n\big)\max_{l}\Big\{\big\|\bm{E}_{l,\cdot}\bm{U}^{(l)}\big\|_{2}+\big\|\bm{E}\big\|\cdot\big\|\bm{U}^{(l)}\big\|_{2,\infty}\Big\}}{\lambda_{r}^{\star}}
%	\label{eq:E-UH-Ustar-UB-234}
%\end{align}
%%


\subsubsection{Step 3: bounding $\|\bm{M}^{\star} (\bm{U}\bm{H}-\bm{U}^{\star} ) \|_{2,\infty}$}


We make the key observation that
%
\begin{align}
\big\|\bm{M}^{\star}\big(\bm{U}\bm{H}-\bm{U}^{\star}\big)\big\|_{2,\infty} & =\big\|\bm{U}^{\star}\bm{\Lambda}^{\star}\bm{U}^{\star\top}\big(\bm{U}\bm{H}-\bm{U}^{\star}\big)\big\|_{2,\infty} \notag\\
 & \leq\big\|\bm{U}^{\star}\big\|_{2,\infty} \|\bm{\Lambda}^{\star}\|\, \big\|\bm{U}^{\star\top}\big(\bm{U}\bm{H}-\bm{U}^{\star}\big)\big\| \notag\\
	& =  \sqrt{\frac{\mu r}{n}} \|\bm{\Lambda}^{\star}\|\, \big\|\bm{U}\bm{U}^{\top} - \bm{U}^{\star} \bm{U}^{\star\top} \big\|^{2}.
	\label{eq:Mstar-UH-Ustar-2-UB-123}
\end{align}
%
Here, the last identity holds true due to the following observation
%
\begin{align*}
\big\|\bm{U}^{\star\top}\big(\bm{U}\bm{H}-\bm{U}^{\star}\big)\big\| & =\big\|\bm{U}^{\star\top}\bm{U}\bm{U}^{\top}\bm{U}^{\star}-\bm{U}^{\star\top}\bm{U}^{\star}\big\|=\big\|\bm{U}^{\star\top}\bm{U}\bm{U}^{\top}\bm{U}^{\star}-\bm{I}\big\|\\
 & =\big\|\bm{Y}(\cos^{2}\bm{\Theta})\bm{Y}^{\top}-\bm{I}\big\|=\big\| \cos^{2}  \bm{\Theta} - \bm{I} \big\|\\
	& =\big\|\sin^{2}\bm{\Theta}\big\|=\big\|\sin\bm{\Theta}\big\|^2,
\end{align*}
%
where we denote by $\bm{X}(\cos \bm{\Theta}) \bm{Y}^{\top}$ the SVD of $\bm{U}^{\top}\bm{U}^{\star}$, with $\bm{X}$ and $\bm{Y}$ being orthonormal matrices and $\bm{\Theta}$ the diagonal matrix consisting of the principal angles between $\bm{U}$ and $\bm{U}^{\star}$ (see the definition in~\eqref{defn:principal-angles}). The above bounds combined with \eqref{eq:distp-U-Ustar-general} indicate that
%
\begin{align}
\big\|\bm{M}^{\star}\big(\bm{U}\bm{H}-\bm{U}^{\star}\big)\big\|_{2,\infty} & \leq
	%\big|\lambda_{1}^{\star}\big|\sqrt{\frac{\mu r}{n}}\,\frac{4c_{2}^{2}\sigma^{2}n}{\big(\lambda_{r}^{\star}\big)^{2}}
	\big|\lambda_{1}^{\star}\big| \sqrt{\frac{\mu r}{n}} \, \big\|\sin\bm{\Theta}\big\|^2
	% = \big|\lambda_{1}^{\star}\big| \sqrt{\frac{\mu r}{n}} \, \big\| \sin \bm{\Theta} \big\|^{2}
	\notag\\
	& \leq \frac{4c_{2}^{2}\kappa\sigma^{2}\sqrt{\mu rn}}{\lambda_{r}^{\star}}
 \eqqcolon \alpha_2.
	\label{eq:Mstar-UH-Ustar-2inf}
\end{align}
%



\subsubsection{Step 4: putting all pieces together}

Combining the bounds \eqref{eq:El-UH-Ustar-UB-Bern123} and \eqref{eq:Mstar-UH-Ustar-2inf} in Steps 2-3 and using the definition of $\mathcal{E}_1$ (see Lemma~\ref{lem:UH-MUstar-diff-decompose}) give
%
\begin{align}
\mathcal{E}_{1} & \leq\frac{2\big\|\bm{E}\big(\bm{U}\bm{H}-\bm{U}^{\star}\big)\big\|_{2,\infty}+2\big\|\bm{M}^{\star}\big(\bm{U}\bm{H}-\bm{U}^{\star}\big)\big\|_{2,\infty}}{\lambda_{r}^{\star}} \nonumber \\
	& \leq\mathcal{E}_{1,1}+\rho_{1}\|\bm{U}\bm{H}-\bm{U}^{\star}\|_{2,\infty},
	\label{eq:E1-UB-general-1}
\end{align}
%
where
%
\begin{equation}
	\mathcal{E}_{1,1}\coloneqq\frac{2\big(\alpha_{0}+\alpha_{1}+\alpha_{2}\big)}{\lambda_{r}^{\star}}
	\quad\text{and}\quad
	\rho_{1}\coloneqq\frac{2c_{4}(\sigma\sqrt{n}+B\log n)}{\lambda_{r}^{\star}}.
	\label{eq:defn-E1-rho1}
\end{equation}
%
This matches precisely the relation hypothesized in \eqref{eq:E1-UB-E11-rho1}. In particular, one has $0<\rho\leq 1/2$
as long as $4c_4(\sigma \sqrt{n}+ B\log n) \leq \lambda_r^{\star}$, which holds whenever $\sigma \sqrt{n \log n}\leq c_{\sigma} \lambda_r^{\star}$ for a sufficiently small $c_{\sigma}>0$ in view of our assumption on $B$ (cf. \eqref{eq:assumption-B-sigma}).


Recall that $\mathcal{E}_{1,1}$, $\mathcal{E}_{2}$, $\mathcal{E}_{3}$, $\alpha_0$, $\alpha_1$, $\alpha_2$ and $\rho_1$ have been controlled in \eqref{eq:defn-E1-rho1}, \eqref{eq:MUstar-E-norm-lambda2-bound}, \eqref{eq:E3-bound-general}, \eqref{eq:defn-alpha-0-general}, \eqref{eq:alpha1-bound-1234}, \eqref{eq:Mstar-UH-Ustar-2inf} and \eqref{eq:defn-E1-rho1}, respectively.
In addition, recall our assumption \eqref{eq:assumption-B-sigma} and suppose that $\sigma{\sqrt{n}}\leq c_{\sigma}\lambda_{r}^{\star}$ for some sufficiently small constant $c_{\sigma}>0$.
With these bounds and assumptions in mind, invoking Lemma~\ref{lem:UH-sequence-bounds-E123} and combining terms immediately conclude the proof.



\begin{remark} We shall also make note of an immediate consequence of the above argument as follows
%
\begin{align}
%\big\|\bm{U}\bm{H}-\bm{M}\bm{U}^{\star}\big(\bm{\Lambda}^{\star}\big)^{-1}\big\|_{2,\infty} & \lesssim\frac{\sigma\kappa\sqrt{\mu r}}{\lambda_{r}^{\star}}+\frac{c_{\mathsf{b}}\sigma^{2}\sqrt{rn}\log n}{\big(\lambda_{r}^{\star}\big)^{2}},
	\big\|\bm{U}\bm{H}-\bm{U}^{\star}\big\|_{2,\infty}&\lesssim\frac{\sigma\kappa\sqrt{\mu r}+\sigma\sqrt{r\log n}}{\lambda_{r}^{\star}},
	\label{eq:UH-MUstar-Lambda-star-2inf}
\end{align}
%
which will prove useful for deriving other important results.
\end{remark}




%and recall the bound \eqref{eq:UH-Ustar-diff-2inf-general}.
%
%
%
%%By rearranging terms, we are left with
%%%
%%\begin{align}
%%\|\bm{U}\bm{H}-\bm{U}^{\star}\|_{2,\infty}
%%	& \leq
%%	\frac{4\big(\alpha_{0}+\alpha_{1} + \alpha_2 \big)}{\lambda_{r}^{\star}}+2\mathcal{E}_{2}+2\mathcal{E}_{3}.
%%	\label{eq:UH-Ustar-general-UB2}
%%\end{align}
%%%
%%Substituting this relation into \eqref{eq:E1-UB-general-1}, we arrive at
%%%
%%\begin{align*}
%%\mathcal{E}_{1} & \leq\frac{2\big(\alpha_{0}+\alpha_{1}+\alpha_{2})+2c_{4}(\sigma\sqrt{n}+B\log n)\Big(\frac{4(\alpha_{0}+\alpha_{1}+\alpha_{2})}{\lambda_{r}^{\star}}+2\mathcal{E}_{2}+3\mathcal{E}_{3}\Big)}{\lambda_{r}^{\star}}\\
%% & \leq\frac{4\big(\alpha_{0}+\alpha_{1}+\alpha_{2})+c_{4}(\sigma\sqrt{n}+B\log n)\Big(4\mathcal{E}_{2}+6\mathcal{E}_{3}\Big)}{\lambda_{r}^{\star}},
%%\end{align*}
%%%
%%provided that $4c_4(\sigma\sqrt{n}+B\log n) \leq \lambda_r^{\star}$.
%
%
%
%Returning to  \eqref{eq:UH-MUstarLambdastar-diff-2inf-general}, we can combine the above bound on $\mathcal{E}_1$ to deduce that
%%
%\begin{align*}
% & \big\|\bm{U}\bm{H}-\bm{M}\bm{U}^{\star}\big(\bm{\Lambda}^{\star}\big)^{-1}\big\|_{2,\infty}\leq\mathcal{E}_{1}+\mathcal{E}_{2}\\
% & \qquad\leq\frac{4\big(\alpha_{0}+\alpha_{1}+\alpha_{2})+6c_{4}(\sigma\sqrt{n}+B\log n)\mathcal{E}_{3}}{\lambda_{r}^{\star}}+2\mathcal{E}_{2},
%\end{align*}
%%
%provided that $4c_4(\sigma\sqrt{n}+B\log n) \leq \lambda_r^{\star}$
%
%
%
%Furthermore,
%%
%\begin{align*}
%\big\|\bm{U}\bm{H}-\bm{U}\mathsf{sgn}(\bm{H})\big\|_{2,\infty} & \leq\big\|\bm{U}\big\|_{2,\infty}\big\|\bm{H}-\mathsf{sgn}(\bm{H})\big\|\leq\frac{4c_{2}^{2}\sigma^{2}n}{(\lambda_{r}^{\star})^{2}}\big\|\bm{U}\bm{H}\big\|_{2,\infty}\\
% & \leq\frac{4c_{2}^{2}\sigma^{2}n}{(\lambda_{r}^{\star})^{2}}\left\{ \big\|\bm{U}\bm{H}-\bm{U}^{\star}\big\|_{2,\infty}+\big\|\bm{U}^{\star}\big\|_{2,\infty}\right\} ,
%\end{align*}
%%
%where the second inequality follows from \eqref{eq:H-property-summary-general} and \eqref{eq:A-2inf-H-Hl-general}. Consequently,
%%
%\begin{align*}
% & \big\|\bm{U}\mathsf{sgn}(\bm{H})-\bm{M}\bm{U}^{\star}\big(\bm{\Lambda}^{\star}\big)^{-1}\big\|_{2,\infty}\\
% & \quad\leq\big\|\bm{U}\bm{H}-\bm{U}\mathsf{sgn}(\bm{H})\big\|_{2,\infty}+\big\|\bm{U}\bm{H}-\bm{M}\bm{U}^{\star}\big(\bm{\Lambda}^{\star}\big)^{-1}\big\|_{2,\infty}\\
% & \quad\leq\frac{4c_{2}^{2}\sigma^{2}n}{(\lambda_{r}^{\star})^{2}}\left\{ \big\|\bm{U}\bm{H}-\bm{U}^{\star}\big\|_{2,\infty}+\sqrt{\frac{\mu r}{n}}\right\} +\big\|\bm{U}\bm{H}-\bm{M}\bm{U}^{\star}\big(\bm{\Lambda}^{\star}\big)^{-1}\big\|_{2,\infty}\\
% & \quad\leq\frac{6\big(\alpha_{0}+\alpha_{1}+\alpha_{2})+8c_{4}(\sigma\sqrt{n}+B\log n)\mathcal{E}_{3}}{\lambda_{r}^{\star}}+3\mathcal{E}_{2}+\frac{4c_{2}^{2}\sigma^{2}\sqrt{\mu nr}}{(\lambda_{r}^{\star})^{2}} .
%\end{align*}
%%
%\yxc{maybe sent earlier}
%


%combine the triangle inequality and the relation \eqref{eq:UH-Ustar-general-UB2} to yield
%%
%\begin{align*}
% & \big\|\bm{U}\bm{H}-\bm{M}\bm{U}^{\star}\big(\bm{\Lambda}^{\star}\big)^{-1}\big\|_{2,\infty}\leq\big\|\bm{U}\bm{H}-\bm{M}^{\star}\bm{U}^{\star}\big(\bm{\Lambda}^{\star}\big)^{-1}+\bm{E}\bm{U}^{\star}\big(\bm{\Lambda}^{\star}\big)^{-1}\big\|_{2,\infty}\\
% & \quad\quad\leq\big\|\bm{U}\bm{H}-\bm{M}^{\star}\bm{U}^{\star}\big(\bm{\Lambda}^{\star}\big)^{-1}\big\|_{2,\infty}+\big\|\bm{E}\bm{U}^{\star}\big(\bm{\Lambda}^{\star}\big)^{-1}\big\|_{2,\infty}\\
% & \quad\quad\leq\big\|\bm{U}\bm{H}-\bm{U}^{\star}\big\|_{2,\infty}+\mathcal{E}_{3}
%	\leq\frac{4\big(\alpha_{0}+\alpha_{1}+\alpha_{2}\big)}{\lambda_{r}^{\star}}+2\mathcal{E}_{2}+3\mathcal{E}_{3}.
%\end{align*}
%%
%


\subsection{Proof of auxiliary lemmas}
\label{sec:proof-auxiliary-lemmas-general}




%\paragraph{Proof of Lemma~\ref{lem:preliminary-result-general-iid}.}
%
%Under Assumption~\ref{assumption-noise-general},  Theorem~\ref{thm:tighter-spectral-normal-ramon} (in particular \eqref{eq:X-spectral-norm-iid-special-ramon})  reveals the existence of some constant $c_2>0$ such that
%%
%\begin{align}
%	\big\|\bm{E}^{(l)} \big\|\leq \|\bm{E}\|\leq c_2 \sigma \sqrt{n}
%	%\qquad\text{and} \qquad
%	% c_2 \sigma \sqrt{n}
%	\qquad 1\leq l \leq n,
%	\label{eq:noise-E-size-general}
%\end{align}
%%
%holds with probability exceeding $1- O(n^{-7})$. Additionally, repeating the argument in the proof of Corollary~\ref{cor:davis-kahan-conclusion-corollary} reveals that
%%
%\begin{subequations}
%\begin{align}
%	\max_{1\leq j\leq r} | \lambda_j  |  &\geq \lambda_r^{\star} - \|\bm{E}\|,  \quad &  \max_{1\leq j\leq r} | \lambda^{(l)}_j  |  &\geq \lambda_r^{\star} - \|\bm{E}^{(l)}\|, &&\\
%	\max_{j: j> r} | \lambda_j   |  &\leq \|\bm{E}\| ,  & \max_{j: j> r}| \lambda_j^{(l)} |  &\leq \|\bm{E}^{(l)}\|,
%\end{align}
%\end{subequations}
%%
%which taken together with \eqref{eq:noise-E-size-general} validates  \eqref{eq:lambda-size-large-general}-\eqref{eq:lambda-size-small-general}.
%
%
%Regarding \eqref{eq:dist-U-Ustar-general} and \eqref{eq:distp-U-Ustar-general}, apply Corollary~\ref{cor:davis-kahan-conclusion-corollary} to reach
%%
%\begin{align}
%	\mathsf{dist}(\bm{U},\bm{U}^{\star}) \leq  \sqrt{2} \|\sin \bm{\Theta} \|
%	\leq\frac{2\|\bm{E}\|}{\lambda_{r}^{\star}}\leq\frac{2c_{2}\sigma\sqrt{n}}{\lambda_{r}^{\star}}	
%	\label{eq:dist-UUstar-sin-Theta-general-UB123}
%\end{align}
%%
%as long as $\|\bm{E}\|\leq c_2 \sigma \sqrt{n}\leq (1-1/\sqrt{2})\lambda_r^{\star}$, where we have used $\lambda_{r+1}^{\star}=0$. The bound on $\mathsf{dist}\big( \bm{U}^{(l)} , \bm{U}^{\star} \big)$ follows from the same argument.
%%
%
%
%When it comes to $\bm{E}\bm{A}$, we proceed by viewing
%its $l$-th row $\bm{E}_{l,\cdot}\bm{A}$ as a sum of independent random
%vectors as follows
%\[
%\bm{E}_{l,\cdot}\bm{A}=\sum\nolimits_{j}E_{l,j}\bm{A}_{j,\cdot}\eqqcolon\sum\nolimits_{j}\bm{z}_{j},
%\]
%which can be controlled by the matrix Bernstein inequality. Specifically,
%it is seen from Assumption~\ref{assumption-noise-general} that
%%
%\begin{align*}
%v & \coloneqq\sum\nolimits_{j}\mathbb{E}\big[\|\bm{z}_{j}\|_{2}^{2}\big]
%	% = \sum_{j}\mathbb{E}\big[E_{l,j}^{2}\big] \mathbb{E}\big[\|\bm{A}_{j,\cdot}\|_{2}^{2}\big]
%	\leq\sigma^{2}\sum\nolimits_{j}\big\|\bm{A}_{j,\cdot}\big\|_{2}^{2}=\sigma^{2}\|\bm{A}\|_{\mathrm{F}}^{2};\\
%L & \coloneqq\max_{j}\|\bm{z}_{j}\|_{2}=\max_{j}|E_{l,j}|\,\big\|\bm{Z}_{j,\cdot}\big\|_{2}\leq B\|\bm{Z}\|_{2,\infty}.
%\end{align*}
%%
%Invoke the matrix Bernstein inequality (cf.~Corollary~\ref{thm:matrix-Bernstein-friendly})
%and take the union bound to demonstrate that: with probability exceeding
%$1-2n^{-6}$,
%%
%\begin{align*}
%	\big\|\bm{E}_{l,\cdot}\bm{A}\big\|_{2} & \leq4\sqrt{v\log n}+6L\log n
%	\leq  4\sigma\sqrt{\log n}\,\|\bm{A}\|_{\mathrm{F}}+(6B\log n)\|\bm{A}\|_{2,\infty}
%\end{align*}
%%
%holds simultaneously for all $1\leq l\leq n$, thus concluding the proof.
%
%
%
%
%\paragraph{Proof of Lemma~\ref{lem:H-property-summary-general}.}
%
%We shall only prove the result for $\bm{H}$; the proof for  $\bm{H}^{(l)}$ follows from identical arguments and is hence omitted.
%
%
%From our discussion in Section~\ref{sec:introduction-distance-principal-angles},
%one can express the SVD of $\bm{H}=\bm{U}^{\top}\bm{U}^{\star}$ as
%$\bm{H}=\bm{X}(\cos\bm{\Theta})\bm{Y}^{\top}$, where the columns
%of $\bm{X}$ (resp.~$\bm{Y}$) are the left (resp.~right) singular
%vectors of $\bm{H}$, and $\bm{\Theta}$ is a diagonal matrix composed
%of the principal angles between $\bm{U}$ and $\bm{U}^{\star}$. In
%light of this  and the definition (\ref{eq:defn-sgn-Z}),
%we can establish \eqref{eq:H-sgnH-dist-general} as follows
%%
%\begin{align}
%\big\|\bm{H}-\mathsf{sgn}(\bm{H})\big\| & =\big\|\bm{X}\big(\cos\bm{\Theta}-\bm{I}\big)\bm{Y}^{\top}\big\|=\big\|\bm{I}-\cos\bm{\Theta}\big\| \notag\\
% & =2\big\|\sin^{2}(\bm{\Theta}/2)\big\|\leq\|\sin\bm{\Theta}\|^{2} \notag\\
% & \leq \frac{2c_{2}^{2}\sigma^{2}n}{(\lambda_{r}^{\star})^{2}},
%	\label{eq:H-sgnH-diff-UB-123}
%\end{align}
%%
%where the middle line holds since $1-\cos\theta=2\sin^{2}(\theta/2)$
%and $\sin(\theta/2)\leq\frac{\sqrt{2}}{2}\sin\theta$ for any $0\leq\theta\leq\pi/2$, and the last line follows from \eqref{eq:distp-U-Ustar-general}.
%
%Coming back to the claim \eqref{eq:H-inv-norm-bound-general}, it suffices to justify that $\sigma_{\min}(\bm{H})\geq 1/2$.
%Recognizing that $\mathsf{sgn}(\bm{H})=\bm{X}\bm{Y}^{\top}$,
%we see that all singular values of $\mathsf{sgn}(\bm{H})$  equal
% 1. Thus, Weyl's inequality together with (\ref{eq:H-sgnH-diff-UB-123})
%gives
%\[
%\sigma_{\min}(\bm{H})\geq\sigma_{\min}\big(\mathsf{sgn}(\bm{H})\big)-\|\bm{H}-\mathsf{sgn}(\bm{H})\|\geq1-\frac{2c_{2}^{2}\sigma^{2}n}{(\lambda_{r}^{\star})^{2}}\geq\frac{1}{2} ,
%\]
%with the proviso that $2c_{2}\sigma\sqrt{n}\leq\lambda_{r}^{\star}$.
%% This concludes the proof.
%
%
%Additionally, regarding $\bm{U}\bm{H}-\bm{U}^{\star}$ we can derive
%%
%\begin{align}
% & \|\bm{U}\bm{H}-\bm{U}^{\star}\|_{\mathrm{F}}
%   = \|\bm{U}\bm{U}^{\top}\bm{U}^{\star}-\bm{U}^{\star}\bm{U}^{\star\top}\bm{U}^{\star}\|_{\mathrm{F}} \nonumber\\
% & \qquad \quad \leq \|\bm{U}\bm{U}^{\top}-\bm{U}^{\star}\bm{U}^{\star\top}\|\,\|\bm{U}^{\star}\|_{\mathrm{F}}
%   \leq\frac{2c_{2}\sigma \sqrt{rn}}{\lambda_{r}^{\star}},
%\end{align}
%%
%where the last line arises from \eqref{eq:distp-U-Ustar-general}, Lemma~\ref{prop:unitary-norm-property} and $\|\bm{U}^{\star}\|_{\mathrm{F}}=\sqrt{r}$.
%
%
%
\paragraph{Proof of Lemma~\ref{lemma:general-noise-bound}.}

Under Assumption~\ref{assumption-noise-general}, Theorem~\ref{thm:tighter-spectral-normal-ramon}
(in particular \eqref{eq:X-spectral-norm-iid-special-ramon}) reveals
the existence of some constant $c_{2}>0$ such that
\[
\max_{l}\big\|\bm{E}^{(l)}\big\|\leq\|\bm{E}\|\leq c_{2}\sigma\sqrt{n},
\]
holds with probability exceeding $1-O(n^{-7})$.

When it comes to $\bm{E}\bm{A}$, we proceed by viewing its $l$-th
row $\bm{E}_{l,\cdot}\bm{A}$ as a sum of independent random vectors
as follows
\[
\bm{E}_{l,\cdot}\bm{A}=\sum\nolimits _{j}E_{l,j}\bm{A}_{j,\cdot}\eqqcolon\sum\nolimits _{j}\bm{z}_{j},
\]
which can be controlled by the matrix Bernstein inequality. Specifically,
it is seen from Assumption~\ref{assumption-noise-general} that
\begin{align*}
v & \coloneqq\sum\nolimits _{j}\mathbb{E}\big[\|\bm{z}_{j}\|_{2}^{2}\big]%=\sum_{j}\mathbb{E}\big[E_{l,j}^{2}\big]\mathbb{E}\big[\|\bm{A}_{j,\cdot}\|_{2}^{2}\big]
\leq\sigma^{2}\sum\nolimits _{j}\big\|\bm{A}_{j,\cdot}\big\|_{2}^{2}=\sigma^{2}\|\bm{A}\|_{\mathrm{F}}^{2};\\
L & \coloneqq\max_{j}\|\bm{z}_{j}\|_{2}=\max_{j}|E_{l,j}|\,\big\|\bm{Z}_{j,\cdot}\big\|_{2}\leq B\|\bm{Z}\|_{2,\infty}.
\end{align*}
Invoke the matrix Bernstein inequality (cf.~Corollary~\ref{thm:matrix-Bernstein-friendly})
and take the union bound to demonstrate that: with probability exceeding
$1-2n^{-6}$,
\begin{align*}
\big\|\bm{E}_{l,\cdot}\bm{A}\big\|_{2} & \leq4\sqrt{v\log n}+6L\log n\leq4\sigma\sqrt{\log n}\,\|\bm{A}\|_{\mathrm{F}}+(6B\log n)\|\bm{A}\|_{2,\infty}
\end{align*}
holds simultaneously for all $1\leq l\leq n$, thus concluding the
proof.

\paragraph{Proof of Lemma~\ref{lem:preliminary-result-general-iid}.}

Lemma~\ref{lemma:general-noise-bound} tells us that
\begin{equation}
\max_{l}\big\|\bm{E}^{(l)}\big\|\leq\|\bm{E}\|\leq c_{2}\sigma\sqrt{n}\label{eq:noise-E-size-general}
\end{equation}
holds with probability exceeding $1-O(n^{-7})$. Repeating the argument
in the proof of Corollary~\ref{cor:davis-kahan-conclusion-corollary}
reveals that \begin{subequations}
\begin{align}
\max_{1\leq j\leq r}|\lambda_{j}| & \geq\lambda_{r}^{\star}-\|\bm{E}\|,\quad & \max_{1\leq j\leq r}|\lambda_{j}^{(l)}| & \geq\lambda_{r}^{\star}-\|\bm{E}^{(l)}\|,\\
\max_{j:j>r}|\lambda_{j}| & \leq\|\bm{E}\|, & \max_{j:j>r}|\lambda_{j}^{(l)}| & \leq\|\bm{E}^{(l)}\|,
\end{align}
\end{subequations} which taken together with \eqref{eq:noise-E-size-general}
validates \eqref{eq:lambda-size-large-general}-\eqref{eq:lambda-size-small-general}.

Regarding \eqref{eq:dist-U-Ustar-general} and \eqref{eq:distp-U-Ustar-general},
apply Corollary~\ref{cor:davis-kahan-conclusion-corollary} to reach
\begin{align}
\mathsf{dist}(\bm{U},\bm{U}^{\star})\leq\sqrt{2}\|\sin\bm{\Theta}\|\leq\frac{2\|\bm{E}\|}{\lambda_{r}^{\star}}\leq\frac{2c_{2}\sigma\sqrt{n}}{\lambda_{r}^{\star}}\label{eq:dist-UUstar-sin-Theta-general-UB123}
\end{align}
as long as $\|\bm{E}\|\leq c_{2}\sigma\sqrt{n}\leq(1-1/\sqrt{2})\lambda_{r}^{\star}$,
where we have used $\lambda_{r+1}^{\star}=0$. The bound on $\mathsf{dist}\big(\bm{U}^{(l)},\bm{U}^{\star}\big)$
follows from the same argument.

Additionally, regarding $\bm{U}\bm{H}-\bm{U}^{\star}$ we can derive
%
\begin{align}
 & \|\bm{U}\bm{H}-\bm{U}^{\star}\|_{\mathrm{F}}=\|\bm{U}\bm{U}^{\top}\bm{U}^{\star}-\bm{U}^{\star}\bm{U}^{\star\top}\bm{U}^{\star}\|_{\mathrm{F}}\nonumber \\
 & \qquad\quad\leq\|\bm{U}\bm{U}^{\top}-\bm{U}^{\star}\bm{U}^{\star\top}\|\,\|\bm{U}^{\star}\|_{\mathrm{F}} \nonumber\\
 &\qquad\quad=\|\sin\bm{\Theta}\|\,\|\bm{U}^{\star}\|_{\mathrm{F}}\leq\frac{2c_{2}\sigma\sqrt{rn}}{\lambda_{r}^{\star}},
\end{align}
where the last line arises from Lemma~\ref{prop:unitary-norm-property},
\eqref{eq:distp-U-Ustar-general}, and $\|\bm{U}^{\star}\|_{\mathrm{F}}=\sqrt{r}$.




\paragraph{Proof of Lemma~\ref{lem:H-property-summary-general}.}

We shall only prove the result for $\bm{H}$; the proof for $\bm{H}^{(l)}$
follows from identical arguments and is hence omitted.

From our discussion in Section~\ref{sec:introduction-distance-principal-angles},
one can express the SVD of $\bm{H}=\bm{U}^{\top}\bm{U}^{\star}$ as
$\bm{H}=\bm{X}(\cos\bm{\Theta})\bm{Y}^{\top}$, where the columns
of $\bm{X}$ (resp.~$\bm{Y}$) are the left (resp.~right) singular
vectors of $\bm{H}$, and $\bm{\Theta}$ is a diagonal matrix composed
of the principal angles between $\bm{U}$ and $\bm{U}^{\star}$. In
light of this and the definition (\ref{eq:defn-sgn-Z}), we can establish
\eqref{eq:H-sgnH-dist-general} as follows
%
\begin{align}
\big\|\bm{H}-\mathsf{sgn}(\bm{H})\big\| & =\big\|\bm{X}\big(\cos\bm{\Theta}-\bm{I}\big)\bm{Y}^{\top}\big\|=\big\|\bm{I}-\cos\bm{\Theta}\big\|\nonumber \\
 & \leq \big\|\bm{I}-\cos^{2}\bm{\Theta}\big\| = \|\sin\bm{\Theta}\|^{2}\nonumber \\
 & \leq\frac{2c_{2}^{2}\sigma^{2}n}{(\lambda_{r}^{\star})^{2}},
	\label{eq:H-sgnH-diff-UB-123}
\end{align}
%
where the middle line holds since $1-\cos\theta \leq 1 - \cos^2 \theta$,
and the last line follows from \eqref{eq:distp-U-Ustar-general}.

Coming back to the claim \eqref{eq:H-inv-norm-bound-general}, it
suffices to justify that $\sigma_{\min}(\bm{H})\geq1/2$. Recognizing
that $\mathsf{sgn}(\bm{H})=\bm{X}\bm{Y}^{\top}$, we see that all
singular values of $\mathsf{sgn}(\bm{H})$ equal 1. Thus, Weyl's inequality
together with (\ref{eq:H-sgnH-diff-UB-123}) gives
\[
\sigma_{\min}(\bm{H})\geq\sigma_{\min}\big(\mathsf{sgn}(\bm{H})\big)-\|\bm{H}-\mathsf{sgn}(\bm{H})\|\geq1-\frac{2c_{2}^{2}\sigma^{2}n}{(\lambda_{r}^{\star})^{2}}\geq\frac{1}{2},
\]
with the proviso that $2c_{2}\sigma\sqrt{n}\leq\lambda_{r}^{\star}$.


\paragraph{Proof of Lemma~\ref{lem:UH-MUstar-diff-decompose}.}
We start by connecting $\bm{U}\bm{H}\bm{\Lambda}^{\star}$ more explicitly with $\bm{M}\bm{U}^{\star}$ as follows (the invertibility of $\bm{\Lambda}$ can be deduced from \eqref{eq:lambda-size-large-general})
%
\begin{align}
	\bm{U}\bm{H}\bm{\Lambda}^{\star} & =\bm{U}\bm{\Lambda}\bm{\Lambda}^{-1}\bm{H}\bm{\Lambda}^{\star}
	%=\bm{M}\bm{U}\bm{\Lambda}^{-1}\bm{U}^{\top}\bm{U}^{\star}
	=\bm{M}\bm{U}\bm{\Lambda}^{-1}\bm{U}^{\top}\bm{U}^{\star}\bm{\Lambda}^{\star},\label{eq:UH-derivation-1}
\end{align}
%
which relies on the definition (\ref{eq:defn-H-Hl-general})
and the eigendecomposition $\bm{M}\bm{U}=\bm{U}\bm{\Lambda}$. In
addition, the eigendecomposition $\bm{M}^{\star}\bm{U}^{\star}=\bm{U}^{\star}\bm{\Lambda}^{\star}$ gives
%
\begin{align}
	\bm{U}^{\top}\bm{U}^{\star}\bm{\Lambda}^{\star}
	&=\bm{U}^{\top}\bm{M}^{\star}\bm{U}^{\star}=\bm{U}^{\top}\bm{M}\bm{U}^{\star}-\bm{U}^{\top}\bm{E}\bm{U}^{\star} \notag\\
	&= \bm{\Lambda}\bm{U}^{\top}\bm{U}^{\star}-\bm{U}^{\top}\bm{E}\bm{U}^{\star},
	\label{eq:HLambda-Lambda-H-connection}
\end{align}
%
which taken collectively with  (\ref{eq:UH-derivation-1}) demonstrates that
%
\begin{align}
\bm{U}\bm{H}\bm{\Lambda}^{\star} & =\bm{M}\bm{U}\bm{\Lambda}^{-1}\bm{\Lambda}\bm{U}^{\top}\bm{U}^{\star}-\bm{M}\bm{U}\bm{\Lambda}^{-1}\bm{U}^{\top}\bm{E}\bm{U}^{\star} \notag\\
 & =\bm{M}\bm{U}\bm{H}-\bm{M}\bm{U}\bm{\Lambda}^{-1}\bm{U}^{\top}\bm{E}\bm{U}^{\star} \notag\\
 & =\bm{M}\bm{U}^{\star}+\bm{M}\big(\bm{U}\bm{H}-\bm{U}^{\star}\big)-\bm{M}\bm{U}\bm{\Lambda}^{-1}\bm{U}^{\top}\bm{E}\bm{U}^{\star}.
\notag
%\label{eq:UH-derivation-2}
\end{align}
%
Consequently, the difference between $\bm{U}\bm{H}\bm{\Lambda}^{\star}$ and $\bm{M}\bm{U}^{\star}$ obeys
%
\begin{align}
 & \big\|\bm{U}\bm{H}\bm{\Lambda}^{\star}-\bm{M}\bm{U}^{\star}\big\|_{2,\infty}
  \leq \big\|\bm{M}\big(\bm{U}\bm{H}-\bm{U}^{\star}\big)\big\|_{2,\infty}  \notag\\
 %& \qquad\qquad\qquad\qquad\qquad\qquad
	& \qquad\qquad\qquad +\big\|\bm{M}\bm{U}\bm{\Lambda}^{-1}\bm{U}^{\top}\bm{E}\bm{U}^{\star}\big\|_{2,\infty}.
\label{eq:UH-MULambda-UB22}
\end{align}
%


%We then upper bound the two terms on the right-hand side of \eqref{eq:UH-MULambda-UB22} separately. Regarding the first term,
%we observe that
%%
%\begin{align}
%  \big\|\bm{M}\big(\bm{U}\bm{H}-\bm{U}^{\star}\big)\big(\bm{\Lambda}^{\star}\big)^{-1}\big\|_{2,\infty}
% & \leq \big\|\bm{M}\big(\bm{U}\bm{H}-\bm{U}^{\star}\big)\big\|_{2,\infty}\big\|\big(\bm{\Lambda}^{\star}\big)^{-1}\big\| \notag\\
% &=\frac{\big\|\bm{M}\big(\bm{U}\bm{H}-\bm{U}^{\star}\big)\big\|_{2,\infty}}{\lambda_{r}^{\star}} ,
%	\label{eq:M-UH-U-diff-UB1234}
%\end{align}
%%
%which follows from the assumption \eqref{eq:assumption-lambda-r-positive}.
%
Regarding the second term in \eqref{eq:UH-MULambda-UB22}, one can deduce that
%\lambda_{r}^{\star}
\begin{align}
 & \big\|\bm{M}\bm{U}\bm{\Lambda}^{-1}\bm{U}^{\top}\bm{E}\bm{U}^{\star}\big\|_{2,\infty} \notag\\
 & \quad \leq\big\|\bm{M}\bm{U}\big\|_{2,\infty}\big\|\bm{\Lambda}^{-1}\big\|\cdot\big\|\bm{U}\big\|\cdot\big\|\bm{E}\big\|\cdot\big\|\bm{U}^{\star}\big\| \notag\\
 & \quad \overset{(\mathrm{i})}{\leq}\frac{2\big\|\bm{M}\bm{U}\big\|_{2,\infty}\big\|\bm{E}\big\|}{\lambda_{r}^{\star}}\overset{(\mathrm{ii})}{\leq}\frac{4\big\|\bm{M}\bm{U}\bm{H}\big\|_{2,\infty}\big\|\bm{E}\big\|}{\lambda_{r}^{\star}} \notag\\
	& \quad \overset{(\mathrm{iii})}{\leq} \frac{4\big\|\bm{M}\big(\bm{U}\bm{H}-\bm{U}^{\star}\big)\big\|_{2,\infty}\big\|\bm{E}\big\|}{ \lambda_{r}^{\star}  }
	+ \frac{4\big\|\bm{M}\bm{U}^{\star}\big\|_{2,\infty}\big\|\bm{E}\big\|}{\lambda_{r}^{\star}} \notag\\
	& \quad \overset{(\mathrm{iv})}{\leq} \big\|\bm{M}\big(\bm{U}\bm{H}-\bm{U}^{\star}\big)\big\|_{2,\infty}
	+ \frac{4\big\|\bm{M}\bm{U}^{\star}\big\|_{2,\infty}\big\|\bm{E}\big\|}{\lambda_{r}^{\star}} .
	\label{eq:M-UH-U-diff-UB5678}
\end{align}
%
Here, (i) follows from the facts $\|\bm{U}\|=\|\bm{U}^{\star}\|=1$ and $|\lambda_r| \geq \lambda_r^{\star} - c_2\sigma\sqrt{n} \geq  \lambda_r^{\star}/2$ (see Lemma~\ref{lem:preliminary-result-general-iid}),
 (ii) holds due to \eqref{eq:A-2inf-H-Hl-general},
% the observation that $\big\|\bm{M}\bm{U}\big\|_{2,\infty}=\big\|\bm{M}\bm{U}\bm{H}\bm{H}^{-1}\big\|_{2,\infty}\leq\big\|\bm{M}\bm{U}\bm{H}\big\|_{2,\infty}\big\|\bm{H}^{-1}\big\|\leq2\big\|\bm{M}\bm{U}\bm{H}\big\|_{2,\infty}$ (given that $\|\bm{H}^{-1}\|\leq 2$ according to \eqref{eq:H-property-summary-general}),
(iii) invokes the triangle inequality,
whereas (iv) holds true provided that $4\|\bm{E}\|\leq \lambda^{\star}_r$.
Combine (\ref{eq:UH-MULambda-UB22}) and \eqref{eq:M-UH-U-diff-UB5678} to reach
\begin{align*}
 & \big\|\bm{U}\bm{H}\bm{\Lambda}^{\star}-\bm{M}\bm{U}^{\star}\big\|_{2,\infty}
  \leq 2\big\|\bm{M}\big(\bm{U}\bm{H}-\bm{U}^{\star}\big)\big\|_{2,\infty}
	+ \frac{4\big\|\bm{M}\bm{U}^{\star}\big\|_{2,\infty}\big\|\bm{E}\big\|}{\lambda_{r}^{\star}},
\end{align*}
which together with the fact that $\|(\bm{\Lambda}^\star)^{-1}\|=1/\lambda^{\star}_{r}$ and the elementary relation $\|\bm{A}\|_{2,\infty}=\|\bm{A}\bm{\Lambda}^\star(\bm{\Lambda}^\star)^{-1}\|_{2,\infty} \leq \|\bm{A}\bm{\Lambda}^\star\|_{2,\infty}\|(\bm{\Lambda}^\star)^{-1}\|$ yields the desired claim \eqref{eq:UH-MUstarLambdastar-diff-2inf-general}.





When it comes to the second claim \eqref{eq:UH-Ustar-diff-2inf-general}, combining \eqref{eq:UH-MUstarLambdastar-diff-2inf-general} with the triangle inequality
%
\begin{align*}
\big\|\bm{U}\bm{H}-\bm{U}^{\star}\big\|_{2,\infty} & = \big\|\bm{U}\bm{H}-\bm{M}^{\star}\bm{U}^{\star}\big(\bm{\Lambda}^{\star}\big)^{-1}\big\|_{2,\infty}\\
 & \leq\big\|\bm{U}\bm{H}-\bm{M}\bm{U}^{\star}\big(\bm{\Lambda}^{\star}\big)^{-1}\big\|_{2,\infty}+\big\|\bm{E}\bm{U}^{\star}\big(\bm{\Lambda}^{\star}\big)^{-1}\big\|_{2,\infty} \\
 & \leq\big\|\bm{U}\bm{H}-\bm{M}\bm{U}^{\star}\big(\bm{\Lambda}^{\star}\big)^{-1}\big\|_{2,\infty}
	+\big\|\bm{E}\bm{U}^{\star}\big\|_{2,\infty} \,/\, \lambda_r^{\star}
\end{align*}
%
immediately establishes the advertised bound. Here the last relation again arises from the elementary inequality $\|\bm{A}\bm{B}\|_{2,\infty}\leq \|\bm{A}\|_{2,\infty}\|\bm{B}\|$.






\paragraph{Proof of Lemma~\ref{lem:UH-sequence-bounds-E123}.}
First of all, taking Condition (\ref{eq:E1-UB-E11-rho1}) collectively with (\ref{eq:UH-Ustar-diff-2inf-general})
and rearranging terms yield (\ref{eq:UH-Ustar-E11-E2-E3}):
%
\begin{equation}
\big\|\bm{U}\bm{H}-\bm{U}^{\star}\big\|_{2,\infty}\leq\frac{1}{1-\rho_{1}}\big(\mathcal{E}_{1,1}+\mathcal{E}_{2}+\mathcal{E}_{3}\big)\leq2\big(\mathcal{E}_{1,1}+\mathcal{E}_{2}+\mathcal{E}_{3}\big),\label{eq:UH-Ustar-E11-E2-E3-1}
\end{equation}
where the last inequality follows from $\rho\leq 1/2$. Substituting (\ref{eq:E1-UB-E11-rho1}) and (\ref{eq:UH-Ustar-E11-E2-E3})
into (\ref{eq:UH-MUstarLambdastar-diff-2inf-general}) then gives \eqref{eq:UH-MUstar-Lambdastar-E11-E2-E3}:
%
\begin{align}
\big\|\bm{U}\bm{H}-\bm{M}\bm{U}^{\star}\big(\bm{\Lambda}^{\star}\big)^{-1}\big\|_{2,\infty} & \leq\mathcal{E}_{1,1}+2\rho_{1}\big(\mathcal{E}_{1,1}+\mathcal{E}_{2}+\mathcal{E}_{3}\big)+\mathcal{E}_{2}\nonumber \\
 & \leq2\mathcal{E}_{1,1}+2\mathcal{E}_{2}+2\rho_{1}\mathcal{E}_{3},\label{eq:UH-MUstar-Lambdastar-E11-E2-E3-1}
\end{align}
where once again we use the assumption that $\rho\leq 1/2$. 
In addition, the following observation connects $\bm{U}\bm{H}$ with
$\bm{U}\mathsf{sgn}(\bm{H})$:
%
\begin{align}
 & \big\|\bm{U}\bm{H}-\bm{U}\mathsf{sgn}(\bm{H})\big\|_{2,\infty}
   \leq \|\bm{U}\|_{2,\infty}\big\|\bm{H}-\mathsf{sgn}(\bm{H})\big\|\nonumber \\
 & \qquad \overset{(\mathrm{i})}{\leq}\frac{2c_{2}^{2}\sigma^{2}n}{(\lambda_{r}^{\star})^{2}}\big\|\bm{U}\big\|_{2,\infty}\overset{(\mathrm{ii})}{\leq}\frac{4c_{2}^{2}\sigma^{2}n}{(\lambda_{r}^{\star})^{2}}\big\|\bm{U}\bm{H}\big\|_{2,\infty}\nonumber \\
 & \qquad \overset{(\mathrm{iii})}{\leq}\frac{4c_{2}^{2}\sigma^{2}n}{(\lambda_{r}^{\star})^{2}}\big\|\bm{U}\bm{H}-\bm{U}^{\star}\big\|_{2,\infty}+\frac{4c_{2}^{2}\sigma^{2}n}{(\lambda_{r}^{\star})^{2}}\sqrt{\frac{\mu r}{n}},
 \label{eq:UH-UsgnH-bound-iterate}
\end{align}
%
where (i) results from (\ref{eq:H-sgnH-dist-general}), (ii) relies on (\ref{eq:A-2inf-H-general}), and (iii) comes from the triangle
inequality $\|\bm{U}\bm{H}\|_{2,\infty}\leq \|\bm{U}\bm{H}-\bm{U}^{\star}\|_{2,\infty} + \|\bm{U}^{\star} \|_{2,\infty}$
and the definition (\ref{eq:defn-Ustar-incoherence}).

The preceding bound
 together with the triangle inequality gives \eqref{eq:UsgnH-Ustar-Lambdastar-E11-E2-E3}:
%
\begin{align*}
\big\|\bm{U}\mathsf{sgn}(\bm{H})-\bm{U}^{\star}\big\|_{2,\infty} & \leq\big\|\bm{U}\bm{H}-\bm{U}^{\star}\big\|_{2,\infty}+\big\|\bm{U}\bm{H}-\bm{U}\mathsf{sgn}(\bm{H})\big\|_{2,\infty}\\
 & \leq\Big(1+\frac{4c_{2}^{2}\sigma^{2}n}{(\lambda_{r}^{\star})^{2}}\Big)\big\|\bm{U}\bm{H}-\bm{U}^{\star}\big\|_{2,\infty}+\frac{4c_{2}^{2}\sigma^{2}n}{(\lambda_{r}^{\star})^{2}}\sqrt{\frac{\mu r}{n}}\\
 & \leq4\big(\mathcal{E}_{1,1}+\mathcal{E}_{2}+\mathcal{E}_{3}\big)+\frac{4c_{2}^{2}\sigma^{2}\sqrt{\mu rn}}{(\lambda_{r}^{\star})^{2}},
\end{align*}
%
where the last inequality holds as long as $4c_{2}^{2}\sigma^{2}n\leq(\lambda_{r}^{\star})^{2}$. Additionally, the inequalities \eqref{eq:UH-MUstar-Lambdastar-E11-E2-E3-1} and \eqref{eq:UH-UsgnH-bound-iterate} further allow us to deduce \eqref{eq:UsgnH-MUstar-Lambdastar-E11-E2-E3}:
%
\begin{align*}
 & \big\|\bm{U}\mathsf{sgn}(\bm{H})-\bm{M}\bm{U}^{\star}\big(\bm{\Lambda}^{\star}\big)^{-1}\big\|_{2,\infty} \\
 & \quad \leq\big\|\bm{U}\bm{H}-\bm{M}\bm{U}^{\star}\big(\bm{\Lambda}^{\star}\big)^{-1}\big\|_{2,\infty}+\big\|\bm{U}\bm{H}-\bm{U}\mathsf{sgn}(\bm{H})\big\|_{2,\infty}\\
 % & \quad\leq2\mathcal{E}_{1,1}+2\mathcal{E}_{2}+2\rho_{1}\mathcal{E}_{3}+\frac{4c_{2}^{2}\sigma^{2}n}{(\lambda_{r}^{\star})^{2}}\big\|\bm{U}\bm{H}-\bm{U}^{\star}\big\|_{2,\infty}+\frac{4c_{2}^{2}\sigma^{2}n}{(\lambda_{r}^{\star})^{2}}\sqrt{\frac{\mu r}{n}}\\
 & \quad\leq2\mathcal{E}_{1,1}+2\mathcal{E}_{2}+2\rho_{1}\mathcal{E}_{3}+\frac{8c_{2}^{2}\sigma^{2}n}{(\lambda_{r}^{\star})^{2}}\big(\mathcal{E}_{1,1}+\mathcal{E}_{2}+\mathcal{E}_{3}\big)+\frac{4c_{2}^{2}\sigma^{2}\sqrt{\mu rn}}{(\lambda_{r}^{\star})^{2}}\\
 & \quad\leq3\mathcal{E}_{1,1}+3\mathcal{E}_{2}+\Big(2\rho_{1}+\frac{8c_{2}^{2}\sigma^{2}n}{(\lambda_{r}^{\star})^{2}}\Big)\mathcal{E}_{3}+\frac{4c_{2}^{2}\sigma^{2}\sqrt{\mu rn}}{(\lambda_{r}^{\star})^{2}} .
\end{align*}
%
Here, the penultimate line combines (\ref{eq:UH-Ustar-E11-E2-E3}), (\ref{eq:UH-MUstar-Lambdastar-E11-E2-E3})
and (\ref{eq:UH-UsgnH-bound-iterate}), while the last inequality relies on the
assumption $8 c_{2}^{2}\sigma^{2}n\leq(\lambda_{r}^{\star})^{2}$.



%
%\begin{align}
% & \big\|\bm{U}\bm{H}-\bm{M}\bm{U}^{\star}\big(\bm{\Lambda}^{\star}\big)^{-1}\big\|_{2,\infty} \notag \\
%	& \qquad \leq\frac{2\big\|\bm{M}\big(\bm{U}\bm{H}-\bm{U}^{\star}\big)\big\|_{2,\infty}}{\lambda_{r}^{\star}}+\frac{4\big\|\bm{M}\bm{U}^{\star}\big\|_{2,\infty}\big\|\bm{E}\big\|}{(\lambda_{r}^{\star})^{2}},
%% & \qquad\leq \frac{2\big\|\bm{E}\big(\bm{U}\bm{H}-\bm{U}^{\star}\big)\big\|_{2,\infty}}{\lambda_{r}^{\star}}  + \frac{2\big\|\bm{M}^{\star}\big(\bm{U}\bm{H}-\bm{U}^{\star}\big)\big\|_{2,\infty}}{\lambda_{r}^{\star}} \notag\\
%%	& \qquad\qquad\quad +\frac{4\Big\{ \sqrt{\frac{\mu r}{n}}\big|\lambda_{1}^{\star}\big|+(4+6c_{1})\sigma\sqrt{r\log n} \Big\} \cdot c_2 \sigma\sqrt{n}}{(\lambda_{r}^{\star})^{2}},
%\label{eq:UB-UH-MULambda-UB3456}
%\end{align}

%
%
%
%
%Substituting \eqref{eq:M-UH-U-diff-UB1234} and \eqref{eq:M-UH-U-diff-UB5678} into \eqref{eq:UH-MULambda-UB22}, we arrive at
%%
%\begin{align}
% & \big\|\bm{U}\bm{H}-\bm{M}\bm{U}^{\star}\big(\bm{\Lambda}^{\star}\big)^{-1}\big\|_{2,\infty}  \leq\frac{2\big\|\bm{M}\big(\bm{U}\bm{H}-\bm{U}^{\star}\big)\big\|_{2,\infty}}{\lambda_{r}^{\star}}+\frac{4\big\|\bm{M}\bm{U}^{\star}\big\|_{2,\infty}\big\|\bm{E}\big\|}{(\lambda_{r}^{\star})^{2}}\nonumber\\
% & \qquad\leq \frac{2\big\|\bm{E}\big(\bm{U}\bm{H}-\bm{U}^{\star}\big)\big\|_{2,\infty}}{\lambda_{r}^{\star}}  + \frac{2\big\|\bm{M}^{\star}\big(\bm{U}\bm{H}-\bm{U}^{\star}\big)\big\|_{2,\infty}}{\lambda_{r}^{\star}} \notag\\
%	& \qquad\qquad\quad +\frac{4\Big\{ \sqrt{\frac{\mu r}{n}}\big|\lambda_{1}^{\star}\big|+(4+6c_{1})\sigma\sqrt{r\log n} \Big\} \cdot c_2 \sigma\sqrt{n}}{(\lambda_{r}^{\star})^{2}},
%\label{eq:UB-UH-MULambda-UB12345}
%\end{align}
%%
%%
%%\begin{align}
%% & \big\|\bm{U}\bm{H}-\bm{M}\bm{U}^{\star}\big(\bm{\Lambda}^{\star}\big)^{-1}\big\|_{2,\infty}\nonumber \\
%% & \leq\frac{2\big\|\bm{M}\big(\bm{U}\bm{H}-\bm{U}^{\star}\big)\big\|_{2,\infty}}{\lambda_{r}^{\star}}+\frac{4\big\|\bm{M}\bm{U}^{\star}\big\|_{2,\infty}\big\|\bm{E}\big\|}{(\lambda_{r}^{\star})^{2}}\nonumber \\
%% &  \leq\frac{2\big\|\bm{M}\big(\bm{U}\bm{H}-\bm{U}^{\star}\big)\big\|_{2,\infty}}{\lambda_{r}^{\star}}
%%	+\frac{4\big\{\sqrt{\frac{\mu r}{n}} \|\bm{\Lambda}^{\star}\|+\big\|\bm{E}\bm{U}^{\star}\big\|_{2,\infty}\big\}\big\|\bm{E}\big\|}{(\lambda_{r}^{\star})^{2}},
%% \label{eq:UB-UH-MULambda-UB12345}
%%\end{align}
%%
%
%





